% =============================================================================
% ΚΕΦΑΛΑΙΟ 6: ΣΥΜΠΕΡΑΣΜΑΤΑ
% =============================================================================
% Αυτό το αρχείο περιέχει το πλήρες Κεφάλαιο 6 της διπλωματικής εργασίας
% "Μοντελοποίηση Αγορών Ηλεκτρικής Ενέργειας με Julia/JuMP"
% 
% Για ενσωμάτωση στο thesis.tex, χρησιμοποιήστε: % =============================================================================
% ΚΕΦΑΛΑΙΟ 6: ΣΥΜΠΕΡΑΣΜΑΤΑ
% =============================================================================
% Αυτό το αρχείο περιέχει το πλήρες Κεφάλαιο 6 της διπλωματικής εργασίας
% "Μοντελοποίηση Αγορών Ηλεκτρικής Ενέργειας με Julia/JuMP"
% 
% Για ενσωμάτωση στο thesis.tex, χρησιμοποιήστε: % =============================================================================
% ΚΕΦΑΛΑΙΟ 6: ΣΥΜΠΕΡΑΣΜΑΤΑ
% =============================================================================
% Αυτό το αρχείο περιέχει το πλήρες Κεφάλαιο 6 της διπλωματικής εργασίας
% "Μοντελοποίηση Αγορών Ηλεκτρικής Ενέργειας με Julia/JuMP"
% 
% Για ενσωμάτωση στο thesis.tex, χρησιμοποιήστε: % =============================================================================
% ΚΕΦΑΛΑΙΟ 6: ΣΥΜΠΕΡΑΣΜΑΤΑ
% =============================================================================
% Αυτό το αρχείο περιέχει το πλήρες Κεφάλαιο 6 της διπλωματικής εργασίας
% "Μοντελοποίηση Αγορών Ηλεκτρικής Ενέργειας με Julia/JuMP"
% 
% Για ενσωμάτωση στο thesis.tex, χρησιμοποιήστε: \input{chapter6_conclusions}
% =============================================================================

\chapter{Συμπεράσματα}
\label{ch:conclusions}

Στο παρόν κεφάλαιο συνοψίζονται τα κύρια ευρήματα και οι συνεισφορές της διπλωματικής εργασίας, αναλύονται οι περιορισμοί της παρούσας υλοποίησης, και προτείνονται κατευθύνσεις για μελλοντική έρευνα και ανάπτυξη.


% =============================================================================
\section{Σύνοψη Εργασίας}
\label{sec:thesis-summary}

Η παρούσα διπλωματική εργασία ασχολήθηκε με τη μοντελοποίηση και προσομοίωση της Προημερήσιας Αγοράς ηλεκτρικής ενέργειας, με έμφαση στον αλγόριθμο EUPHEMIA που χρησιμοποιείται για την εκκαθάριση των ευρωπαϊκών αγορών. Αναπτύχθηκε ένα ολοκληρωμένο σύστημα λογισμικού που καλύπτει όλα τα στάδια της διαδικασίας, από τη συλλογή δεδομένων έως τον υπολογισμό τιμών εκκαθάρισης.

Η εργασία ξεκίνησε με την ανάλυση του πλαισίου λειτουργίας των αγορών ηλεκτρικής ενέργειας στην Ευρώπη, εξετάζοντας την ιστορική εξέλιξη από τα κρατικά μονοπώλια στις απελευθερωμένες αγορές. Παρουσιάστηκε η δομή της Προημερήσιας Αγοράς και ο ρόλος του αλγορίθμου EUPHEMIA στην ενοποίηση των ευρωπαϊκών αγορών μέσω του μηχανισμού Single Day-Ahead Coupling (SDAC).

Στη συνέχεια, αναπτύχθηκε το θεωρητικό υπόβαθρο που απαιτείται για την κατανόηση του προβλήματος. Παρουσιάστηκαν οι βασικές έννοιες του μαθηματικού προγραμματισμού, με έμφαση στα προβλήματα μεικτού ακέραιου προγραμματισμού (MIP) και στα προβλήματα συμπληρωματικότητας (MPCC). Διατυπώθηκε το πρόβλημα Unit Commitment και αναλύθηκε ο αλγόριθμος EUPHEMIA με βάση τη δημόσια τεκμηρίωση.

Η αρχιτεκτονική του συστήματος σχεδιάστηκε με γνώμονα την επεκτασιμότητα και τη συντηρησιμότητα. Αναπτύχθηκε αγωγός δεδομένων (data pipeline) σε Python για τη συλλογή δεδομένων από την πλατφόρμα διαφάνειας του ENTSO-E, με αποθήκευση σε βάση δεδομένων PostgreSQL. Η κύρια εφαρμογή υλοποιήθηκε στη γλώσσα Julia με χρήση του πλαισίου JuMP για τη μοντελοποίηση των προβλημάτων βελτιστοποίησης.

Η υλοποίηση περιλαμβάνει τρία κύρια συστατικά. Πρώτον, τον επιλυτή Unit Commitment που βελτιστοποιεί τη δέσμευση και την κατανομή των μονάδων παραγωγής λαμβάνοντας υπόψη τους τεχνικούς περιορισμούς κάθε τύπου καυσίμου. Δεύτερον, τη μηχανή στρατηγικής προσφορών που μετατρέπει τα αποτελέσματα του Unit Commitment σε εντολές αγοράς με διάφορες στρατηγικές. Τρίτον, τον αλγόριθμο εκκαθάρισης αγοράς βασισμένο στη μέθοδο MPCC, που υπολογίζει τις τιμές εκκαθάρισης και τους βαθμούς αποδοχής των εντολών.

Τέλος, το σύστημα αξιολογήθηκε μέσω σύγκρισης των υπολογισμένων τιμών με τις πραγματικές τιμές της αγοράς. Το βιβλίο εντολών οριακού κόστους, τροφοδοτούμενο με τους κύριους θεμελιώδεις παράγοντες της αγοράς (ημερήσιες τιμές φυσικού αερίου TTF και δικαιωμάτων CO$_2$, αξία νερού συναρτήσει της πληρότητας των ταμιευτήρων, αυτοπρογραμματισμός μονάδων βάσης, παρατηρούμενες καθαρές εισαγωγές), αναπαράγει τις τιμές της ελληνικής ζώνης με συντελεστή συσχέτισης $r = 0{,}84$ σε ανάλυση δεκαπενταλέπτου επί συνεχούς τετραμήνου, με πρακτικά μηδενική συστηματική απόκλιση και διάμεσο ενδοημερήσιο συντελεστή συσχέτισης 0{,}90 (Κεφάλαιο 5).


% =============================================================================
\section{Κύριες Συνεισφορές}
\label{sec:contributions}

Η παρούσα εργασία παρουσιάζει τις ακόλουθες συνεισφορές στους τομείς της μοντελοποίησης αγορών ενέργειας και της εφαρμοσμένης βελτιστοποίησης.

\subsection{Υλοποίηση Ανοιχτού Κώδικα}
\label{subsec:open-source-contribution}

Η πρώτη και κυριότερη συνεισφορά είναι η ανάπτυξη μιας πλήρως λειτουργικής υλοποίησης ανοιχτού κώδικα για την προσομοίωση της Προημερήσιας Αγοράς. Ενώ ο αλγόριθμος EUPHEMIA είναι δημόσια τεκμηριωμένος, οι υπάρχουσες υλοποιήσεις είναι είτε ιδιόκτητες είτε περιορισμένης λειτουργικότητας. Εξ όσων είμαστε σε θέση να γνωρίζουμε, πρόκειται για το πρώτο σύστημα ανοιχτού κώδικα που μοντελοποιεί την αγορά αυτή με τεκμηριωμένη, ποσοτικά επικυρωμένη ακρίβεια σε πραγματικά δεδομένα και σε συνεχή περίοδο λειτουργίας, με πλήρως αναπαραγώγιμη μεθοδολογία --- από τα ακατέργαστα δημόσια δεδομένα έως τις τιμές εκκαθάρισης. Το σύστημα που αναπτύχθηκε στην παρούσα εργασία διατίθεται ελεύθερα στη διεύθυνση \url{https://github.com/philokalia-ai/energy-markets} και μπορεί να χρησιμοποιηθεί για εκπαιδευτικούς και ερευνητικούς σκοπούς.

Η υλοποίηση περιλαμβάνει όλα τα απαραίτητα συστατικά για την εκτέλεση προσομοιώσεων, από τη συλλογή δεδομένων έως την αποθήκευση αποτελεσμάτων. Ο κώδικας είναι καλά τεκμηριωμένος και δομημένος με τρόπο που διευκολύνει την επέκταση και την προσαρμογή σε διαφορετικές ανάγκες.

\subsection{Ολοκληρωμένος Αγωγός Δεδομένων}
\label{subsec:data-pipeline-contribution}

Η δεύτερη συνεισφορά είναι η ανάπτυξη ενός πλήρους αγωγού δεδομένων για τη συλλογή και επεξεργασία δεδομένων από την πλατφόρμα διαφάνειας του ENTSO-E. Ο αγωγός υποστηρίζει σταδιακές ενημερώσεις (incremental updates), αποτελεσματική μαζική φόρτωση δεδομένων, και αυτοματοποιημένη εκτέλεση μέσω GitHub Actions.

Τα δεδομένα που συλλέγονται περιλαμβάνουν πληροφορίες για τις μονάδες παραγωγής, το ιστορικό φορτίο, τις προβλέψεις παραγωγής από ΑΠΕ, τις διαθέσιμες ικανότητες μεταφοράς, και τις μη διαθεσιμότητες μονάδων. Ο σχεδιασμός του σχήματος της βάσης δεδομένων επιτρέπει την αποτελεσματική εκτέλεση ερωτημάτων και τη διατήρηση της ιστορικότητας των δεδομένων.

\subsection{Πολυζωνική Εκκαθάριση Αγοράς}
\label{subsec:multizone-contribution}

Η τρίτη συνεισφορά είναι η υλοποίηση πολυζωνικής εκκαθάρισης αγοράς με περιορισμούς διασυνοριακής μεταφοράς. Το σύστημα μοντελοποιεί τη διαθέσιμη ικανότητα μεταφοράς (Available Transfer Capacity - ATC) μεταξύ ζωνών προσφοράς και υπολογίζει ταυτόχρονα τις τιμές σε όλες τις συνδεδεμένες ζώνες.

Η προσέγγιση αυτή επιτρέπει την ανάλυση της επίδρασης της διασυνδεσιμότητας στις τιμές και τον εντοπισμό περιπτώσεων συμφόρησης που οδηγούν σε διαφοροποίηση τιμών μεταξύ ζωνών. Τα αποτελέσματα περιλαμβάνουν όχι μόνο τις τιμές ανά ζώνη αλλά και τις ροές μεταφοράς μεταξύ ζωνών.

\subsection{Ανταγωνιστικό Αντιπαράδειγμα με Επικυρωμένη Ακρίβεια}
\label{subsec:counterfactual-contribution}

Σημαντική συνεισφορά αποτελεί η μεθοδολογία του βιβλίου εντολών οριακού κόστους (Κεφάλαιο 4), η οποία κατασκευάζει ένα \emph{ανταγωνιστικό αντιπαράδειγμα} της αγοράς: τις τιμές που θα προέκυπταν εάν κάθε συμμετέχων προσέφερε στο βραχυχρόνιο οριακό του κόστος συν την ορθολογική διαχείριση των τεχνικών του περιορισμών. Η προσέγγιση επιτυγχάνει συσχέτιση 0{,}84 σε ανάλυση δεκαπενταλέπτου χωρίς καμία στατιστική προσαρμογή εκ των υστέρων, διατηρώντας πλήρη ερμηνευσιμότητα: κάθε υπολογισμένη τιμή αντιστοιχεί σε συγκεκριμένη οριακή εντολή συγκεκριμένης τεχνολογίας. Πέραν της προσομοίωσης, η ιδιότητα αυτή καθιστά το σύστημα εργαλείο εποπτείας αγοράς: συστηματικές αποκλίσεις των πραγματικών τιμών \emph{πάνω} από το αντιπαράδειγμα, ιδίως όταν το ενδοημερήσιο σχήμα εξηγείται πλήρως, συνιστούν ποσοτικό μέτρο της απόκλισης της πραγματικής συμπεριφοράς προσφορών από την ανταγωνιστική (Ενότητα 5.7).

\subsection{Επίδειξη του Οικοσυστήματος Julia/JuMP}
\label{subsec:julia-contribution}

Η τέταρτη συνεισφορά είναι η επίδειξη της καταλληλότητας του οικοσυστήματος Julia/JuMP για προβλήματα βελτιστοποίησης στον τομέα της ενέργειας. Η Julia συνδυάζει την εκφραστικότητα γλωσσών υψηλού επιπέδου με την υπολογιστική απόδοση γλωσσών χαμηλού επιπέδου, ενώ το JuMP παρέχει μια διαισθητική διεπαφή για τη μοντελοποίηση προβλημάτων βελτιστοποίησης.

Η εργασία επιδεικνύει τη χρήση του JuMP για τη διατύπωση πολύπλοκων προβλημάτων βελτιστοποίησης, τη σύνδεση με διάφορους επιλυτές (HiGHS, Gurobi), και την αλληλεπίδραση με βάσεις δεδομένων. Ο κώδικας μπορεί να χρησιμεύσει ως παράδειγμα για παρόμοιες εφαρμογές.

\subsection{Μοντελοποίηση Ειδικών Παραμέτρων ανά Τύπο Καυσίμου}
\label{subsec:fuel-type-contribution}

Η πέμπτη συνεισφορά είναι η λεπτομερής μοντελοποίηση των τεχνικών χαρακτηριστικών των μονάδων παραγωγής ανάλογα με τον τύπο καυσίμου. Το σύστημα διαφοροποιεί τους χρόνους εκκίνησης, τους ρυθμούς μεταβολής, και τους ελάχιστους χρόνους λειτουργίας μεταξύ διαφορετικών τύπων καυσίμου όπως λιγνίτης, φυσικό αέριο, και υδροηλεκτρικά.

Η διαφοροποίηση αυτή οδηγεί σε ρεαλιστικότερα αποτελέσματα Unit Commitment και κατ' επέκταση σε ακριβέστερες προσφορές στην αγορά.


% =============================================================================
\section{Περιορισμοί}
\label{sec:limitations}

Παρά τις συνεισφορές που αναφέρθηκαν, η παρούσα υλοποίηση υπόκειται σε ορισμένους περιορισμούς που πρέπει να ληφθούν υπόψη κατά την ερμηνεία των αποτελεσμάτων.

\subsection{Μονοπερίοδη Βελτιστοποίηση}
\label{subsec:single-period-limitation}

Το σύστημα εκτελεί βελτιστοποίηση για μία ημέρα κάθε φορά, χωρίς να λαμβάνει υπόψη την αλληλεξάρτηση μεταξύ διαδοχικών ημερών. Στην πραγματικότητα, οι αποφάσεις δέσμευσης μονάδων επηρεάζονται από τις συνθήκες των προηγούμενων ημερών (αρχικές καταστάσεις) και τις προβλέψεις για τις επόμενες ημέρες (κυλιόμενος ορίζοντας).

Η απουσία κυλιόμενου ορίζοντα μπορεί να οδηγήσει σε μη βέλτιστες αποφάσεις, ιδιαίτερα για μονάδες με μεγάλους χρόνους εκκίνησης που πρέπει να ξεκινήσουν νωρίς για να είναι διαθέσιμες σε περιόδους υψηλής ζήτησης.

\subsection{Απλοποιημένο Μοντέλο Δικτύου}
\label{subsec:network-limitation}

Το μοντέλο δικτύου βασίζεται αποκλειστικά στη διαθέσιμη ικανότητα μεταφοράς (ATC) μεταξύ ζωνών προσφοράς, χωρίς να μοντελοποιεί τους νόμους της ροής ισχύος (Kirchhoff). Η προσέγγιση ATC είναι η τυπική για τις αγορές ηλεκτρικής ενέργειας, αλλά δεν αποτυπώνει με ακρίβεια τη φυσική συμπεριφορά του δικτύου.

Επιπλέον, δεν μοντελοποιούνται οι απώλειες μεταφοράς, οι οποίες μπορεί να είναι σημαντικές για μεγάλες αποστάσεις μεταφοράς.

\subsection{Περιορισμένοι Τύποι Εντολών}
\label{subsec:order-types-limitation}

Η υλοποίηση υποστηρίζει κυρίως απλές ωριαίες εντολές (Simple Orders), ενώ ο πραγματικός αλγόριθμος EUPHEMIA υποστηρίζει πολύ πιο σύνθετους τύπους εντολών. Οι εντολές μπλοκ (Block Orders) που επιτρέπουν την υποβολή προσφοράς για πολλαπλές συνεχόμενες ώρες υπό συνθήκη ελάχιστου εσόδου δεν υλοποιούνται πλήρως.

Επίσης, δεν υποστηρίζονται οι συνδεδεμένες εντολές μπλοκ (Linked Block Orders), οι αποκλειστικές εντολές μπλοκ (Exclusive Block Orders), οι εντολές ελάχιστου εισοδήματος (MIC Orders), και άλλοι προηγμένοι τύποι.

\subsection{Εκτίμηση Οικονομικών Παραμέτρων}
\label{subsec:economic-parameters-limitation}

Τα οικονομικά χαρακτηριστικά των μονάδων παραγωγής (οριακά κόστη, κόστη εκκίνησης, κόστη no-load) δεν δημοσιεύονται από το ENTSO-E και πρέπει να εκτιμηθούν με βάση τον τύπο καυσίμου και δημόσια διαθέσιμες πληροφορίες για τις τιμές καυσίμων και τους ρύπους CO$_2$.

Η αβεβαιότητα στην εκτίμηση αυτών των παραμέτρων αποτελεί σημαντική πηγή σφάλματος στα αποτελέσματα της προσομοίωσης.

\subsection{Μη Μοντελοποίηση Στρατηγικής Συμπεριφοράς}
\label{subsec:strategic-behavior-limitation}

Το σύστημα υποθέτει ότι οι συμμετέχοντες στην αγορά υποβάλλουν προσφορές με βάση το πραγματικό τους κόστος (ή με σταθερό συντελεστή markup). Στην πραγματικότητα, οι συμμετέχοντες μπορεί να ακολουθούν στρατηγικές υποβολής προσφορών που λαμβάνουν υπόψη τη θέση τους στην αγορά και τις αναμενόμενες ενέργειες των ανταγωνιστών.

Η επιλογή αυτή είναι εν μέρει σκόπιμη: καθιστά το σύστημα ανταγωνιστικό αντιπαράδειγμα, οπότε οι αποκλίσεις των πραγματικών τιμών από τις υπολογισμένες δεν αποτελούν αποκλειστικά σφάλμα μοντελοποίησης αλλά και μετρήσιμη ένδειξη μη ανταγωνιστικής συμπεριφοράς προσφορών. Περιορίζει ωστόσο την ικανότητα αναπαραγωγής των πραγματικών τιμών σε ημέρες στενότητας, όπου η στρατηγική συμπεριφορά είναι εντονότερη.


% =============================================================================
\section{Μελλοντική Εργασία}
\label{sec:future-work}

Με βάση την εμπειρία που αποκτήθηκε κατά την ανάπτυξη του συστήματος και τους περιορισμούς που εντοπίστηκαν, προτείνονται οι ακόλουθες κατευθύνσεις για μελλοντική έρευνα και ανάπτυξη.

\subsection{Συστηματοποίηση της Εποπτείας Αγοράς}
\label{subsec:market-monitoring-future}

Άμεση προέκταση της παρούσας εργασίας είναι η συστηματοποίηση της ανάλυσης καταλοίπων ως μεθοδολογίας εποπτείας αγοράς. Η ανάλυση σε επίπεδο μονάδας που παρουσιάστηκε ενδεικτικά στο Κεφάλαιο 5 (σύγκριση πραγματικής κατανομής με την αναμενόμενη από το ανταγωνιστικό αντιπαράδειγμα, εντοπισμός αδήλωτα μη διαθέσιμης ισχύος, δείκτης εναπομένοντος προμηθευτή) μπορεί να αυτοματοποιηθεί σε όλες τις ζώνες και ημέρες, με παλινδρόμηση των συναγόμενων περιθωρίων προσφορών σε εκ των προτέρων παρατηρήσιμα σήματα (περιθώριο επάρκειας, πρόβλεψη ΑΠΕ, πληρότητα ταμιευτήρων, ανακοινωθείσες μη διαθεσιμότητες) και έλεγχο συσχέτισης μεταξύ εταιρειών πέραν των κοινών σημάτων.

\subsection{Επέκταση στη Σύζευξη Βάσει Ροών}
\label{subsec:flow-based-future}

Η ανάλυση ανέδειξε ότι μετά τον Ιούνιο του 2022 τα σύνορα της περιοχής Core εκκαθαρίζονται με σύζευξη βάσει ροών (flow-based market coupling) και δεν αναπαρίστανται από δεδομένα ATC. Η υλοποίηση της flow-based προσέγγισης (πολύγωνα εφικτότητας από PTDF και περιθώρια κρίσιμων στοιχείων δικτύου) θα επέτρεπε την ενδογενή αναπαράσταση ζωνών όπως η Ουγγαρία, οι οποίες σήμερα αποδίδονται ατελώς ως δέκτες τιμών.

\subsection{Επέκταση σε Ενδοημερήσια Αγορά}
\label{subsec:intraday-extension}

Η πρώτη κατεύθυνση είναι η επέκταση του συστήματος για την προσομοίωση της Ενδοημερήσιας Αγοράς (Intraday Market). Η Ενδοημερήσια Αγορά λειτουργεί συνεχώς μεταξύ της λήξης της Προημερήσιας και της έναρξης της ημέρας παράδοσης, επιτρέποντας στους συμμετέχοντες να προσαρμόζουν τις θέσεις τους με βάση ενημερωμένες προβλέψεις.

Η μοντελοποίηση της Ενδοημερήσιας Αγοράς απαιτεί διαφορετικούς μηχανισμούς εκκαθάρισης (συνεχής διαπραγμάτευση vs δημοπρασία) και θα παρείχε πληρέστερη εικόνα της λειτουργίας της αγοράς.

\subsection{Προσομοίωση Αγοράς Εξισορρόπησης}
\label{subsec:balancing-simulation}

Η δεύτερη κατεύθυνση είναι η προσομοίωση της Αγοράς Εξισορρόπησης (Balancing Market). Η αγορά αυτή λειτουργεί σε πραγματικό χρόνο για την αντιμετώπιση αποκλίσεων μεταξύ παραγωγής και κατανάλωσης, και αποτελεί κρίσιμο συστατικό της λειτουργίας του συστήματος ηλεκτρικής ενέργειας.

Η μοντελοποίηση της Αγοράς Εξισορρόπησης θα απαιτήσει τη συμπερίληψη αβεβαιοτήτων στις προβλέψεις ζήτησης και παραγωγής ΑΠΕ, καθώς και τη μοντελοποίηση των αποθεμάτων ισχύος (reserves).

\subsection{Ενσωμάτωση Μηχανικής Μάθησης}
\label{subsec:ml-integration}

Η τρίτη κατεύθυνση είναι η ενσωμάτωση τεχνικών μηχανικής μάθησης για τη βελτίωση διαφόρων συνιστωσών του συστήματος. Συγκεκριμένα, η μηχανική μάθηση μπορεί να χρησιμοποιηθεί για την πρόβλεψη τιμών ηλεκτρικής ενέργειας, την εκτίμηση οικονομικών παραμέτρων μονάδων παραγωγής, και την πρόβλεψη διαθεσιμότητας μονάδων βάσει ιστορικών δεδομένων μη διαθεσιμότητας.

Η χρήση υβριδικών μοντέλων που συνδυάζουν τη φυσική γνώση (μοντέλα βελτιστοποίησης) με τη μηχανική μάθηση (data-driven approaches) αποτελεί πολλά υποσχόμενη κατεύθυνση.

\subsection{Υλοποίηση Σύνθετων Τύπων Εντολών}
\label{subsec:complex-orders}

Η τέταρτη κατεύθυνση είναι η πλήρης υλοποίηση των σύνθετων τύπων εντολών που υποστηρίζει ο αλγόριθμος EUPHEMIA. Οι εντολές μπλοκ με συνθήκη ελάχιστου εσόδου, οι συνδεδεμένες εντολές, και οι εντολές ελάχιστου εισοδήματος είναι απαραίτητες για τη ρεαλιστική αναπαράσταση της συμπεριφοράς θερμικών μονάδων στην αγορά.

Η υλοποίηση αυτών των τύπων εντολών θα απαιτήσει την επέκταση του μοντέλου MPCC με επιπλέον μεταβλητές και περιορισμούς.

\subsection{Βελτιστοποίηση Υπολογιστικής Απόδοσης}
\label{subsec:performance-optimization}

Η πέμπτη κατεύθυνση είναι η βελτιστοποίηση της υπολογιστικής απόδοσης του συστήματος για τη διαχείριση μεγαλύτερων προβλημάτων. Η παραλληλοποίηση της επίλυσης για πολλαπλές ημερομηνίες, η χρήση τεχνικών warm-start για διαδοχικές επιλύσεις, και η εξερεύνηση εναλλακτικών διατυπώσεων του προβλήματος αποτελούν πιθανές κατευθύνσεις.

Επιπλέον, η ενσωμάτωση του συστήματος σε περιβάλλον cloud θα επέτρεπε την κλιμάκωση σε μεγαλύτερα προβλήματα και την εκτέλεση εκτεταμένων αναλύσεων ευαισθησίας.

\subsection{Επέκταση σε Περιφερειακές Αγορές}
\label{subsec:regional-extension}

Η έκτη κατεύθυνση είναι η επέκταση της ανάλυσης σε περισσότερες ζώνες προσφοράς και η μελέτη περιφερειακών αγορών πέρα από τη Νοτιοανατολική Ευρώπη. Η σύγκριση των αποτελεσμάτων μεταξύ διαφορετικών περιφερειών θα παρείχε πληροφορίες για τις διαφορές στη δομή και λειτουργία των αγορών.

Η επέκταση αυτή θα απαιτήσει την αντιμετώπιση διαφορών στη διαθεσιμότητα και ποιότητα δεδομένων μεταξύ ζωνών.


% =============================================================================
\section{Επίλογος}
\label{sec:epilogue}

Η παρούσα διπλωματική εργασία επιχείρησε να συμβάλει στην κατανόηση και μοντελοποίηση των αγορών ηλεκτρικής ενέργειας μέσω της ανάπτυξης ενός ολοκληρωμένου συστήματος προσομοίωσης. Η εργασία κινείται στη διασταύρωση πολλών επιστημονικών πεδίων: της επιχειρησιακής έρευνας και βελτιστοποίησης, της οικονομικής θεωρίας των αγορών, της μηχανικής δεδομένων, και της ανάπτυξης λογισμικού.

Η σημασία της κατανόησης των μηχανισμών της αγοράς ηλεκτρικής ενέργειας αυξάνεται συνεχώς καθώς η ενεργειακή μετάβαση επιταχύνεται. Η διείσδυση των ανανεώσιμων πηγών ενέργειας, η ηλεκτροκίνηση, και η αποκέντρωση της παραγωγής δημιουργούν νέες προκλήσεις για τη λειτουργία των αγορών. Τα εργαλεία προσομοίωσης όπως αυτό που αναπτύχθηκε στην παρούσα εργασία μπορούν να συμβάλουν στην ανάλυση των επιπτώσεων αυτών των αλλαγών και στον σχεδιασμό πολιτικών.

Η επιλογή της Julia και του JuMP για την υλοποίηση αποδείχθηκε επιτυχής, επιτρέποντας τη διατύπωση των προβλημάτων βελτιστοποίησης με τρόπο που μοιάζει με τη μαθηματική διατύπωση, διατηρώντας παράλληλα υψηλή υπολογιστική απόδοση. Η εμπειρία αυτή επιβεβαιώνει ότι το οικοσύστημα Julia είναι ώριμο για χρήση σε εφαρμογές βελτιστοποίησης στον τομέα της ενέργειας.

Κλείνοντας, η εργασία αυτή αποτελεί μια βάση πάνω στην οποία μπορεί να χτιστεί περαιτέρω έρευνα και ανάπτυξη. Η διάθεση του κώδικα ως ανοιχτού λογισμικού έχει ως στόχο να διευκολύνει αυτή τη διαδικασία και να συμβάλει στην ευρύτερη κοινότητα που ασχολείται με τη μοντελοποίηση ενεργειακών συστημάτων.

% =============================================================================

\chapter{Συμπεράσματα}
\label{ch:conclusions}

Στο παρόν κεφάλαιο συνοψίζονται τα κύρια ευρήματα και οι συνεισφορές της διπλωματικής εργασίας, αναλύονται οι περιορισμοί της παρούσας υλοποίησης, και προτείνονται κατευθύνσεις για μελλοντική έρευνα και ανάπτυξη.


% =============================================================================
\section{Σύνοψη Εργασίας}
\label{sec:thesis-summary}

Η παρούσα διπλωματική εργασία ασχολήθηκε με τη μοντελοποίηση και προσομοίωση της Προημερήσιας Αγοράς ηλεκτρικής ενέργειας, με έμφαση στον αλγόριθμο EUPHEMIA που χρησιμοποιείται για την εκκαθάριση των ευρωπαϊκών αγορών. Αναπτύχθηκε ένα ολοκληρωμένο σύστημα λογισμικού που καλύπτει όλα τα στάδια της διαδικασίας, από τη συλλογή δεδομένων έως τον υπολογισμό τιμών εκκαθάρισης.

Η εργασία ξεκίνησε με την ανάλυση του πλαισίου λειτουργίας των αγορών ηλεκτρικής ενέργειας στην Ευρώπη, εξετάζοντας την ιστορική εξέλιξη από τα κρατικά μονοπώλια στις απελευθερωμένες αγορές. Παρουσιάστηκε η δομή της Προημερήσιας Αγοράς και ο ρόλος του αλγορίθμου EUPHEMIA στην ενοποίηση των ευρωπαϊκών αγορών μέσω του μηχανισμού Single Day-Ahead Coupling (SDAC).

Στη συνέχεια, αναπτύχθηκε το θεωρητικό υπόβαθρο που απαιτείται για την κατανόηση του προβλήματος. Παρουσιάστηκαν οι βασικές έννοιες του μαθηματικού προγραμματισμού, με έμφαση στα προβλήματα μεικτού ακέραιου προγραμματισμού (MIP) και στα προβλήματα συμπληρωματικότητας (MPCC). Διατυπώθηκε το πρόβλημα Unit Commitment και αναλύθηκε ο αλγόριθμος EUPHEMIA με βάση τη δημόσια τεκμηρίωση.

Η αρχιτεκτονική του συστήματος σχεδιάστηκε με γνώμονα την επεκτασιμότητα και τη συντηρησιμότητα. Αναπτύχθηκε αγωγός δεδομένων (data pipeline) σε Python για τη συλλογή δεδομένων από την πλατφόρμα διαφάνειας του ENTSO-E, με αποθήκευση σε βάση δεδομένων PostgreSQL. Η κύρια εφαρμογή υλοποιήθηκε στη γλώσσα Julia με χρήση του πλαισίου JuMP για τη μοντελοποίηση των προβλημάτων βελτιστοποίησης.

Η υλοποίηση περιλαμβάνει τρία κύρια συστατικά. Πρώτον, τον επιλυτή Unit Commitment που βελτιστοποιεί τη δέσμευση και την κατανομή των μονάδων παραγωγής λαμβάνοντας υπόψη τους τεχνικούς περιορισμούς κάθε τύπου καυσίμου. Δεύτερον, τη μηχανή στρατηγικής προσφορών που μετατρέπει τα αποτελέσματα του Unit Commitment σε εντολές αγοράς με διάφορες στρατηγικές. Τρίτον, τον αλγόριθμο εκκαθάρισης αγοράς βασισμένο στη μέθοδο MPCC, που υπολογίζει τις τιμές εκκαθάρισης και τους βαθμούς αποδοχής των εντολών.

Τέλος, το σύστημα αξιολογήθηκε μέσω σύγκρισης των υπολογισμένων τιμών με τις πραγματικές τιμές της αγοράς. Το βιβλίο εντολών οριακού κόστους, τροφοδοτούμενο με τους κύριους θεμελιώδεις παράγοντες της αγοράς (ημερήσιες τιμές φυσικού αερίου TTF και δικαιωμάτων CO$_2$, αξία νερού συναρτήσει της πληρότητας των ταμιευτήρων, αυτοπρογραμματισμός μονάδων βάσης, παρατηρούμενες καθαρές εισαγωγές), αναπαράγει τις τιμές της ελληνικής ζώνης με συντελεστή συσχέτισης $r = 0{,}84$ σε ανάλυση δεκαπενταλέπτου επί συνεχούς τετραμήνου, με πρακτικά μηδενική συστηματική απόκλιση και διάμεσο ενδοημερήσιο συντελεστή συσχέτισης 0{,}90 (Κεφάλαιο 5).


% =============================================================================
\section{Κύριες Συνεισφορές}
\label{sec:contributions}

Η παρούσα εργασία παρουσιάζει τις ακόλουθες συνεισφορές στους τομείς της μοντελοποίησης αγορών ενέργειας και της εφαρμοσμένης βελτιστοποίησης.

\subsection{Υλοποίηση Ανοιχτού Κώδικα}
\label{subsec:open-source-contribution}

Η πρώτη και κυριότερη συνεισφορά είναι η ανάπτυξη μιας πλήρως λειτουργικής υλοποίησης ανοιχτού κώδικα για την προσομοίωση της Προημερήσιας Αγοράς. Ενώ ο αλγόριθμος EUPHEMIA είναι δημόσια τεκμηριωμένος, οι υπάρχουσες υλοποιήσεις είναι είτε ιδιόκτητες είτε περιορισμένης λειτουργικότητας. Εξ όσων είμαστε σε θέση να γνωρίζουμε, πρόκειται για το πρώτο σύστημα ανοιχτού κώδικα που μοντελοποιεί την αγορά αυτή με τεκμηριωμένη, ποσοτικά επικυρωμένη ακρίβεια σε πραγματικά δεδομένα και σε συνεχή περίοδο λειτουργίας, με πλήρως αναπαραγώγιμη μεθοδολογία --- από τα ακατέργαστα δημόσια δεδομένα έως τις τιμές εκκαθάρισης. Το σύστημα που αναπτύχθηκε στην παρούσα εργασία διατίθεται ελεύθερα στη διεύθυνση \url{https://github.com/philokalia-ai/energy-markets} και μπορεί να χρησιμοποιηθεί για εκπαιδευτικούς και ερευνητικούς σκοπούς.

Η υλοποίηση περιλαμβάνει όλα τα απαραίτητα συστατικά για την εκτέλεση προσομοιώσεων, από τη συλλογή δεδομένων έως την αποθήκευση αποτελεσμάτων. Ο κώδικας είναι καλά τεκμηριωμένος και δομημένος με τρόπο που διευκολύνει την επέκταση και την προσαρμογή σε διαφορετικές ανάγκες.

\subsection{Ολοκληρωμένος Αγωγός Δεδομένων}
\label{subsec:data-pipeline-contribution}

Η δεύτερη συνεισφορά είναι η ανάπτυξη ενός πλήρους αγωγού δεδομένων για τη συλλογή και επεξεργασία δεδομένων από την πλατφόρμα διαφάνειας του ENTSO-E. Ο αγωγός υποστηρίζει σταδιακές ενημερώσεις (incremental updates), αποτελεσματική μαζική φόρτωση δεδομένων, και αυτοματοποιημένη εκτέλεση μέσω GitHub Actions.

Τα δεδομένα που συλλέγονται περιλαμβάνουν πληροφορίες για τις μονάδες παραγωγής, το ιστορικό φορτίο, τις προβλέψεις παραγωγής από ΑΠΕ, τις διαθέσιμες ικανότητες μεταφοράς, και τις μη διαθεσιμότητες μονάδων. Ο σχεδιασμός του σχήματος της βάσης δεδομένων επιτρέπει την αποτελεσματική εκτέλεση ερωτημάτων και τη διατήρηση της ιστορικότητας των δεδομένων.

\subsection{Πολυζωνική Εκκαθάριση Αγοράς}
\label{subsec:multizone-contribution}

Η τρίτη συνεισφορά είναι η υλοποίηση πολυζωνικής εκκαθάρισης αγοράς με περιορισμούς διασυνοριακής μεταφοράς. Το σύστημα μοντελοποιεί τη διαθέσιμη ικανότητα μεταφοράς (Available Transfer Capacity - ATC) μεταξύ ζωνών προσφοράς και υπολογίζει ταυτόχρονα τις τιμές σε όλες τις συνδεδεμένες ζώνες.

Η προσέγγιση αυτή επιτρέπει την ανάλυση της επίδρασης της διασυνδεσιμότητας στις τιμές και τον εντοπισμό περιπτώσεων συμφόρησης που οδηγούν σε διαφοροποίηση τιμών μεταξύ ζωνών. Τα αποτελέσματα περιλαμβάνουν όχι μόνο τις τιμές ανά ζώνη αλλά και τις ροές μεταφοράς μεταξύ ζωνών.

\subsection{Ανταγωνιστικό Αντιπαράδειγμα με Επικυρωμένη Ακρίβεια}
\label{subsec:counterfactual-contribution}

Σημαντική συνεισφορά αποτελεί η μεθοδολογία του βιβλίου εντολών οριακού κόστους (Κεφάλαιο 4), η οποία κατασκευάζει ένα \emph{ανταγωνιστικό αντιπαράδειγμα} της αγοράς: τις τιμές που θα προέκυπταν εάν κάθε συμμετέχων προσέφερε στο βραχυχρόνιο οριακό του κόστος συν την ορθολογική διαχείριση των τεχνικών του περιορισμών. Η προσέγγιση επιτυγχάνει συσχέτιση 0{,}84 σε ανάλυση δεκαπενταλέπτου χωρίς καμία στατιστική προσαρμογή εκ των υστέρων, διατηρώντας πλήρη ερμηνευσιμότητα: κάθε υπολογισμένη τιμή αντιστοιχεί σε συγκεκριμένη οριακή εντολή συγκεκριμένης τεχνολογίας. Πέραν της προσομοίωσης, η ιδιότητα αυτή καθιστά το σύστημα εργαλείο εποπτείας αγοράς: συστηματικές αποκλίσεις των πραγματικών τιμών \emph{πάνω} από το αντιπαράδειγμα, ιδίως όταν το ενδοημερήσιο σχήμα εξηγείται πλήρως, συνιστούν ποσοτικό μέτρο της απόκλισης της πραγματικής συμπεριφοράς προσφορών από την ανταγωνιστική (Ενότητα 5.7).

\subsection{Επίδειξη του Οικοσυστήματος Julia/JuMP}
\label{subsec:julia-contribution}

Η τέταρτη συνεισφορά είναι η επίδειξη της καταλληλότητας του οικοσυστήματος Julia/JuMP για προβλήματα βελτιστοποίησης στον τομέα της ενέργειας. Η Julia συνδυάζει την εκφραστικότητα γλωσσών υψηλού επιπέδου με την υπολογιστική απόδοση γλωσσών χαμηλού επιπέδου, ενώ το JuMP παρέχει μια διαισθητική διεπαφή για τη μοντελοποίηση προβλημάτων βελτιστοποίησης.

Η εργασία επιδεικνύει τη χρήση του JuMP για τη διατύπωση πολύπλοκων προβλημάτων βελτιστοποίησης, τη σύνδεση με διάφορους επιλυτές (HiGHS, Gurobi), και την αλληλεπίδραση με βάσεις δεδομένων. Ο κώδικας μπορεί να χρησιμεύσει ως παράδειγμα για παρόμοιες εφαρμογές.

\subsection{Μοντελοποίηση Ειδικών Παραμέτρων ανά Τύπο Καυσίμου}
\label{subsec:fuel-type-contribution}

Η πέμπτη συνεισφορά είναι η λεπτομερής μοντελοποίηση των τεχνικών χαρακτηριστικών των μονάδων παραγωγής ανάλογα με τον τύπο καυσίμου. Το σύστημα διαφοροποιεί τους χρόνους εκκίνησης, τους ρυθμούς μεταβολής, και τους ελάχιστους χρόνους λειτουργίας μεταξύ διαφορετικών τύπων καυσίμου όπως λιγνίτης, φυσικό αέριο, και υδροηλεκτρικά.

Η διαφοροποίηση αυτή οδηγεί σε ρεαλιστικότερα αποτελέσματα Unit Commitment και κατ' επέκταση σε ακριβέστερες προσφορές στην αγορά.


% =============================================================================
\section{Περιορισμοί}
\label{sec:limitations}

Παρά τις συνεισφορές που αναφέρθηκαν, η παρούσα υλοποίηση υπόκειται σε ορισμένους περιορισμούς που πρέπει να ληφθούν υπόψη κατά την ερμηνεία των αποτελεσμάτων.

\subsection{Μονοπερίοδη Βελτιστοποίηση}
\label{subsec:single-period-limitation}

Το σύστημα εκτελεί βελτιστοποίηση για μία ημέρα κάθε φορά, χωρίς να λαμβάνει υπόψη την αλληλεξάρτηση μεταξύ διαδοχικών ημερών. Στην πραγματικότητα, οι αποφάσεις δέσμευσης μονάδων επηρεάζονται από τις συνθήκες των προηγούμενων ημερών (αρχικές καταστάσεις) και τις προβλέψεις για τις επόμενες ημέρες (κυλιόμενος ορίζοντας).

Η απουσία κυλιόμενου ορίζοντα μπορεί να οδηγήσει σε μη βέλτιστες αποφάσεις, ιδιαίτερα για μονάδες με μεγάλους χρόνους εκκίνησης που πρέπει να ξεκινήσουν νωρίς για να είναι διαθέσιμες σε περιόδους υψηλής ζήτησης.

\subsection{Απλοποιημένο Μοντέλο Δικτύου}
\label{subsec:network-limitation}

Το μοντέλο δικτύου βασίζεται αποκλειστικά στη διαθέσιμη ικανότητα μεταφοράς (ATC) μεταξύ ζωνών προσφοράς, χωρίς να μοντελοποιεί τους νόμους της ροής ισχύος (Kirchhoff). Η προσέγγιση ATC είναι η τυπική για τις αγορές ηλεκτρικής ενέργειας, αλλά δεν αποτυπώνει με ακρίβεια τη φυσική συμπεριφορά του δικτύου.

Επιπλέον, δεν μοντελοποιούνται οι απώλειες μεταφοράς, οι οποίες μπορεί να είναι σημαντικές για μεγάλες αποστάσεις μεταφοράς.

\subsection{Περιορισμένοι Τύποι Εντολών}
\label{subsec:order-types-limitation}

Η υλοποίηση υποστηρίζει κυρίως απλές ωριαίες εντολές (Simple Orders), ενώ ο πραγματικός αλγόριθμος EUPHEMIA υποστηρίζει πολύ πιο σύνθετους τύπους εντολών. Οι εντολές μπλοκ (Block Orders) που επιτρέπουν την υποβολή προσφοράς για πολλαπλές συνεχόμενες ώρες υπό συνθήκη ελάχιστου εσόδου δεν υλοποιούνται πλήρως.

Επίσης, δεν υποστηρίζονται οι συνδεδεμένες εντολές μπλοκ (Linked Block Orders), οι αποκλειστικές εντολές μπλοκ (Exclusive Block Orders), οι εντολές ελάχιστου εισοδήματος (MIC Orders), και άλλοι προηγμένοι τύποι.

\subsection{Εκτίμηση Οικονομικών Παραμέτρων}
\label{subsec:economic-parameters-limitation}

Τα οικονομικά χαρακτηριστικά των μονάδων παραγωγής (οριακά κόστη, κόστη εκκίνησης, κόστη no-load) δεν δημοσιεύονται από το ENTSO-E και πρέπει να εκτιμηθούν με βάση τον τύπο καυσίμου και δημόσια διαθέσιμες πληροφορίες για τις τιμές καυσίμων και τους ρύπους CO$_2$.

Η αβεβαιότητα στην εκτίμηση αυτών των παραμέτρων αποτελεί σημαντική πηγή σφάλματος στα αποτελέσματα της προσομοίωσης.

\subsection{Μη Μοντελοποίηση Στρατηγικής Συμπεριφοράς}
\label{subsec:strategic-behavior-limitation}

Το σύστημα υποθέτει ότι οι συμμετέχοντες στην αγορά υποβάλλουν προσφορές με βάση το πραγματικό τους κόστος (ή με σταθερό συντελεστή markup). Στην πραγματικότητα, οι συμμετέχοντες μπορεί να ακολουθούν στρατηγικές υποβολής προσφορών που λαμβάνουν υπόψη τη θέση τους στην αγορά και τις αναμενόμενες ενέργειες των ανταγωνιστών.

Η επιλογή αυτή είναι εν μέρει σκόπιμη: καθιστά το σύστημα ανταγωνιστικό αντιπαράδειγμα, οπότε οι αποκλίσεις των πραγματικών τιμών από τις υπολογισμένες δεν αποτελούν αποκλειστικά σφάλμα μοντελοποίησης αλλά και μετρήσιμη ένδειξη μη ανταγωνιστικής συμπεριφοράς προσφορών. Περιορίζει ωστόσο την ικανότητα αναπαραγωγής των πραγματικών τιμών σε ημέρες στενότητας, όπου η στρατηγική συμπεριφορά είναι εντονότερη.


% =============================================================================
\section{Μελλοντική Εργασία}
\label{sec:future-work}

Με βάση την εμπειρία που αποκτήθηκε κατά την ανάπτυξη του συστήματος και τους περιορισμούς που εντοπίστηκαν, προτείνονται οι ακόλουθες κατευθύνσεις για μελλοντική έρευνα και ανάπτυξη.

\subsection{Συστηματοποίηση της Εποπτείας Αγοράς}
\label{subsec:market-monitoring-future}

Άμεση προέκταση της παρούσας εργασίας είναι η συστηματοποίηση της ανάλυσης καταλοίπων ως μεθοδολογίας εποπτείας αγοράς. Η ανάλυση σε επίπεδο μονάδας που παρουσιάστηκε ενδεικτικά στο Κεφάλαιο 5 (σύγκριση πραγματικής κατανομής με την αναμενόμενη από το ανταγωνιστικό αντιπαράδειγμα, εντοπισμός αδήλωτα μη διαθέσιμης ισχύος, δείκτης εναπομένοντος προμηθευτή) μπορεί να αυτοματοποιηθεί σε όλες τις ζώνες και ημέρες, με παλινδρόμηση των συναγόμενων περιθωρίων προσφορών σε εκ των προτέρων παρατηρήσιμα σήματα (περιθώριο επάρκειας, πρόβλεψη ΑΠΕ, πληρότητα ταμιευτήρων, ανακοινωθείσες μη διαθεσιμότητες) και έλεγχο συσχέτισης μεταξύ εταιρειών πέραν των κοινών σημάτων.

\subsection{Επέκταση στη Σύζευξη Βάσει Ροών}
\label{subsec:flow-based-future}

Η ανάλυση ανέδειξε ότι μετά τον Ιούνιο του 2022 τα σύνορα της περιοχής Core εκκαθαρίζονται με σύζευξη βάσει ροών (flow-based market coupling) και δεν αναπαρίστανται από δεδομένα ATC. Η υλοποίηση της flow-based προσέγγισης (πολύγωνα εφικτότητας από PTDF και περιθώρια κρίσιμων στοιχείων δικτύου) θα επέτρεπε την ενδογενή αναπαράσταση ζωνών όπως η Ουγγαρία, οι οποίες σήμερα αποδίδονται ατελώς ως δέκτες τιμών.

\subsection{Επέκταση σε Ενδοημερήσια Αγορά}
\label{subsec:intraday-extension}

Η πρώτη κατεύθυνση είναι η επέκταση του συστήματος για την προσομοίωση της Ενδοημερήσιας Αγοράς (Intraday Market). Η Ενδοημερήσια Αγορά λειτουργεί συνεχώς μεταξύ της λήξης της Προημερήσιας και της έναρξης της ημέρας παράδοσης, επιτρέποντας στους συμμετέχοντες να προσαρμόζουν τις θέσεις τους με βάση ενημερωμένες προβλέψεις.

Η μοντελοποίηση της Ενδοημερήσιας Αγοράς απαιτεί διαφορετικούς μηχανισμούς εκκαθάρισης (συνεχής διαπραγμάτευση vs δημοπρασία) και θα παρείχε πληρέστερη εικόνα της λειτουργίας της αγοράς.

\subsection{Προσομοίωση Αγοράς Εξισορρόπησης}
\label{subsec:balancing-simulation}

Η δεύτερη κατεύθυνση είναι η προσομοίωση της Αγοράς Εξισορρόπησης (Balancing Market). Η αγορά αυτή λειτουργεί σε πραγματικό χρόνο για την αντιμετώπιση αποκλίσεων μεταξύ παραγωγής και κατανάλωσης, και αποτελεί κρίσιμο συστατικό της λειτουργίας του συστήματος ηλεκτρικής ενέργειας.

Η μοντελοποίηση της Αγοράς Εξισορρόπησης θα απαιτήσει τη συμπερίληψη αβεβαιοτήτων στις προβλέψεις ζήτησης και παραγωγής ΑΠΕ, καθώς και τη μοντελοποίηση των αποθεμάτων ισχύος (reserves).

\subsection{Ενσωμάτωση Μηχανικής Μάθησης}
\label{subsec:ml-integration}

Η τρίτη κατεύθυνση είναι η ενσωμάτωση τεχνικών μηχανικής μάθησης για τη βελτίωση διαφόρων συνιστωσών του συστήματος. Συγκεκριμένα, η μηχανική μάθηση μπορεί να χρησιμοποιηθεί για την πρόβλεψη τιμών ηλεκτρικής ενέργειας, την εκτίμηση οικονομικών παραμέτρων μονάδων παραγωγής, και την πρόβλεψη διαθεσιμότητας μονάδων βάσει ιστορικών δεδομένων μη διαθεσιμότητας.

Η χρήση υβριδικών μοντέλων που συνδυάζουν τη φυσική γνώση (μοντέλα βελτιστοποίησης) με τη μηχανική μάθηση (data-driven approaches) αποτελεί πολλά υποσχόμενη κατεύθυνση.

\subsection{Υλοποίηση Σύνθετων Τύπων Εντολών}
\label{subsec:complex-orders}

Η τέταρτη κατεύθυνση είναι η πλήρης υλοποίηση των σύνθετων τύπων εντολών που υποστηρίζει ο αλγόριθμος EUPHEMIA. Οι εντολές μπλοκ με συνθήκη ελάχιστου εσόδου, οι συνδεδεμένες εντολές, και οι εντολές ελάχιστου εισοδήματος είναι απαραίτητες για τη ρεαλιστική αναπαράσταση της συμπεριφοράς θερμικών μονάδων στην αγορά.

Η υλοποίηση αυτών των τύπων εντολών θα απαιτήσει την επέκταση του μοντέλου MPCC με επιπλέον μεταβλητές και περιορισμούς.

\subsection{Βελτιστοποίηση Υπολογιστικής Απόδοσης}
\label{subsec:performance-optimization}

Η πέμπτη κατεύθυνση είναι η βελτιστοποίηση της υπολογιστικής απόδοσης του συστήματος για τη διαχείριση μεγαλύτερων προβλημάτων. Η παραλληλοποίηση της επίλυσης για πολλαπλές ημερομηνίες, η χρήση τεχνικών warm-start για διαδοχικές επιλύσεις, και η εξερεύνηση εναλλακτικών διατυπώσεων του προβλήματος αποτελούν πιθανές κατευθύνσεις.

Επιπλέον, η ενσωμάτωση του συστήματος σε περιβάλλον cloud θα επέτρεπε την κλιμάκωση σε μεγαλύτερα προβλήματα και την εκτέλεση εκτεταμένων αναλύσεων ευαισθησίας.

\subsection{Επέκταση σε Περιφερειακές Αγορές}
\label{subsec:regional-extension}

Η έκτη κατεύθυνση είναι η επέκταση της ανάλυσης σε περισσότερες ζώνες προσφοράς και η μελέτη περιφερειακών αγορών πέρα από τη Νοτιοανατολική Ευρώπη. Η σύγκριση των αποτελεσμάτων μεταξύ διαφορετικών περιφερειών θα παρείχε πληροφορίες για τις διαφορές στη δομή και λειτουργία των αγορών.

Η επέκταση αυτή θα απαιτήσει την αντιμετώπιση διαφορών στη διαθεσιμότητα και ποιότητα δεδομένων μεταξύ ζωνών.


% =============================================================================
\section{Επίλογος}
\label{sec:epilogue}

Η παρούσα διπλωματική εργασία επιχείρησε να συμβάλει στην κατανόηση και μοντελοποίηση των αγορών ηλεκτρικής ενέργειας μέσω της ανάπτυξης ενός ολοκληρωμένου συστήματος προσομοίωσης. Η εργασία κινείται στη διασταύρωση πολλών επιστημονικών πεδίων: της επιχειρησιακής έρευνας και βελτιστοποίησης, της οικονομικής θεωρίας των αγορών, της μηχανικής δεδομένων, και της ανάπτυξης λογισμικού.

Η σημασία της κατανόησης των μηχανισμών της αγοράς ηλεκτρικής ενέργειας αυξάνεται συνεχώς καθώς η ενεργειακή μετάβαση επιταχύνεται. Η διείσδυση των ανανεώσιμων πηγών ενέργειας, η ηλεκτροκίνηση, και η αποκέντρωση της παραγωγής δημιουργούν νέες προκλήσεις για τη λειτουργία των αγορών. Τα εργαλεία προσομοίωσης όπως αυτό που αναπτύχθηκε στην παρούσα εργασία μπορούν να συμβάλουν στην ανάλυση των επιπτώσεων αυτών των αλλαγών και στον σχεδιασμό πολιτικών.

Η επιλογή της Julia και του JuMP για την υλοποίηση αποδείχθηκε επιτυχής, επιτρέποντας τη διατύπωση των προβλημάτων βελτιστοποίησης με τρόπο που μοιάζει με τη μαθηματική διατύπωση, διατηρώντας παράλληλα υψηλή υπολογιστική απόδοση. Η εμπειρία αυτή επιβεβαιώνει ότι το οικοσύστημα Julia είναι ώριμο για χρήση σε εφαρμογές βελτιστοποίησης στον τομέα της ενέργειας.

Κλείνοντας, η εργασία αυτή αποτελεί μια βάση πάνω στην οποία μπορεί να χτιστεί περαιτέρω έρευνα και ανάπτυξη. Η διάθεση του κώδικα ως ανοιχτού λογισμικού έχει ως στόχο να διευκολύνει αυτή τη διαδικασία και να συμβάλει στην ευρύτερη κοινότητα που ασχολείται με τη μοντελοποίηση ενεργειακών συστημάτων.

% =============================================================================

\chapter{Συμπεράσματα}
\label{ch:conclusions}

Στο παρόν κεφάλαιο συνοψίζονται τα κύρια ευρήματα και οι συνεισφορές της διπλωματικής εργασίας, αναλύονται οι περιορισμοί της παρούσας υλοποίησης, και προτείνονται κατευθύνσεις για μελλοντική έρευνα και ανάπτυξη.


% =============================================================================
\section{Σύνοψη Εργασίας}
\label{sec:thesis-summary}

Η παρούσα διπλωματική εργασία ασχολήθηκε με τη μοντελοποίηση και προσομοίωση της Προημερήσιας Αγοράς ηλεκτρικής ενέργειας, με έμφαση στον αλγόριθμο EUPHEMIA που χρησιμοποιείται για την εκκαθάριση των ευρωπαϊκών αγορών. Αναπτύχθηκε ένα ολοκληρωμένο σύστημα λογισμικού που καλύπτει όλα τα στάδια της διαδικασίας, από τη συλλογή δεδομένων έως τον υπολογισμό τιμών εκκαθάρισης.

Η εργασία ξεκίνησε με την ανάλυση του πλαισίου λειτουργίας των αγορών ηλεκτρικής ενέργειας στην Ευρώπη, εξετάζοντας την ιστορική εξέλιξη από τα κρατικά μονοπώλια στις απελευθερωμένες αγορές. Παρουσιάστηκε η δομή της Προημερήσιας Αγοράς και ο ρόλος του αλγορίθμου EUPHEMIA στην ενοποίηση των ευρωπαϊκών αγορών μέσω του μηχανισμού Single Day-Ahead Coupling (SDAC).

Στη συνέχεια, αναπτύχθηκε το θεωρητικό υπόβαθρο που απαιτείται για την κατανόηση του προβλήματος. Παρουσιάστηκαν οι βασικές έννοιες του μαθηματικού προγραμματισμού, με έμφαση στα προβλήματα μεικτού ακέραιου προγραμματισμού (MIP) και στα προβλήματα συμπληρωματικότητας (MPCC). Διατυπώθηκε το πρόβλημα Unit Commitment και αναλύθηκε ο αλγόριθμος EUPHEMIA με βάση τη δημόσια τεκμηρίωση.

Η αρχιτεκτονική του συστήματος σχεδιάστηκε με γνώμονα την επεκτασιμότητα και τη συντηρησιμότητα. Αναπτύχθηκε αγωγός δεδομένων (data pipeline) σε Python για τη συλλογή δεδομένων από την πλατφόρμα διαφάνειας του ENTSO-E, με αποθήκευση σε βάση δεδομένων PostgreSQL. Η κύρια εφαρμογή υλοποιήθηκε στη γλώσσα Julia με χρήση του πλαισίου JuMP για τη μοντελοποίηση των προβλημάτων βελτιστοποίησης.

Η υλοποίηση περιλαμβάνει τρία κύρια συστατικά. Πρώτον, τον επιλυτή Unit Commitment που βελτιστοποιεί τη δέσμευση και την κατανομή των μονάδων παραγωγής λαμβάνοντας υπόψη τους τεχνικούς περιορισμούς κάθε τύπου καυσίμου. Δεύτερον, τη μηχανή στρατηγικής προσφορών που μετατρέπει τα αποτελέσματα του Unit Commitment σε εντολές αγοράς με διάφορες στρατηγικές. Τρίτον, τον αλγόριθμο εκκαθάρισης αγοράς βασισμένο στη μέθοδο MPCC, που υπολογίζει τις τιμές εκκαθάρισης και τους βαθμούς αποδοχής των εντολών.

Τέλος, το σύστημα αξιολογήθηκε μέσω σύγκρισης των υπολογισμένων τιμών με τις πραγματικές τιμές της αγοράς. Το βιβλίο εντολών οριακού κόστους, τροφοδοτούμενο με τους κύριους θεμελιώδεις παράγοντες της αγοράς (ημερήσιες τιμές φυσικού αερίου TTF και δικαιωμάτων CO$_2$, αξία νερού συναρτήσει της πληρότητας των ταμιευτήρων, αυτοπρογραμματισμός μονάδων βάσης, παρατηρούμενες καθαρές εισαγωγές), αναπαράγει τις τιμές της ελληνικής ζώνης με συντελεστή συσχέτισης $r = 0{,}84$ σε ανάλυση δεκαπενταλέπτου επί συνεχούς τετραμήνου, με πρακτικά μηδενική συστηματική απόκλιση και διάμεσο ενδοημερήσιο συντελεστή συσχέτισης 0{,}90 (Κεφάλαιο 5).


% =============================================================================
\section{Κύριες Συνεισφορές}
\label{sec:contributions}

Η παρούσα εργασία παρουσιάζει τις ακόλουθες συνεισφορές στους τομείς της μοντελοποίησης αγορών ενέργειας και της εφαρμοσμένης βελτιστοποίησης.

\subsection{Υλοποίηση Ανοιχτού Κώδικα}
\label{subsec:open-source-contribution}

Η πρώτη και κυριότερη συνεισφορά είναι η ανάπτυξη μιας πλήρως λειτουργικής υλοποίησης ανοιχτού κώδικα για την προσομοίωση της Προημερήσιας Αγοράς. Ενώ ο αλγόριθμος EUPHEMIA είναι δημόσια τεκμηριωμένος, οι υπάρχουσες υλοποιήσεις είναι είτε ιδιόκτητες είτε περιορισμένης λειτουργικότητας. Εξ όσων είμαστε σε θέση να γνωρίζουμε, πρόκειται για το πρώτο σύστημα ανοιχτού κώδικα που μοντελοποιεί την αγορά αυτή με τεκμηριωμένη, ποσοτικά επικυρωμένη ακρίβεια σε πραγματικά δεδομένα και σε συνεχή περίοδο λειτουργίας, με πλήρως αναπαραγώγιμη μεθοδολογία --- από τα ακατέργαστα δημόσια δεδομένα έως τις τιμές εκκαθάρισης. Το σύστημα που αναπτύχθηκε στην παρούσα εργασία διατίθεται ελεύθερα στη διεύθυνση \url{https://github.com/philokalia-ai/energy-markets} και μπορεί να χρησιμοποιηθεί για εκπαιδευτικούς και ερευνητικούς σκοπούς.

Η υλοποίηση περιλαμβάνει όλα τα απαραίτητα συστατικά για την εκτέλεση προσομοιώσεων, από τη συλλογή δεδομένων έως την αποθήκευση αποτελεσμάτων. Ο κώδικας είναι καλά τεκμηριωμένος και δομημένος με τρόπο που διευκολύνει την επέκταση και την προσαρμογή σε διαφορετικές ανάγκες.

\subsection{Ολοκληρωμένος Αγωγός Δεδομένων}
\label{subsec:data-pipeline-contribution}

Η δεύτερη συνεισφορά είναι η ανάπτυξη ενός πλήρους αγωγού δεδομένων για τη συλλογή και επεξεργασία δεδομένων από την πλατφόρμα διαφάνειας του ENTSO-E. Ο αγωγός υποστηρίζει σταδιακές ενημερώσεις (incremental updates), αποτελεσματική μαζική φόρτωση δεδομένων, και αυτοματοποιημένη εκτέλεση μέσω GitHub Actions.

Τα δεδομένα που συλλέγονται περιλαμβάνουν πληροφορίες για τις μονάδες παραγωγής, το ιστορικό φορτίο, τις προβλέψεις παραγωγής από ΑΠΕ, τις διαθέσιμες ικανότητες μεταφοράς, και τις μη διαθεσιμότητες μονάδων. Ο σχεδιασμός του σχήματος της βάσης δεδομένων επιτρέπει την αποτελεσματική εκτέλεση ερωτημάτων και τη διατήρηση της ιστορικότητας των δεδομένων.

\subsection{Πολυζωνική Εκκαθάριση Αγοράς}
\label{subsec:multizone-contribution}

Η τρίτη συνεισφορά είναι η υλοποίηση πολυζωνικής εκκαθάρισης αγοράς με περιορισμούς διασυνοριακής μεταφοράς. Το σύστημα μοντελοποιεί τη διαθέσιμη ικανότητα μεταφοράς (Available Transfer Capacity - ATC) μεταξύ ζωνών προσφοράς και υπολογίζει ταυτόχρονα τις τιμές σε όλες τις συνδεδεμένες ζώνες.

Η προσέγγιση αυτή επιτρέπει την ανάλυση της επίδρασης της διασυνδεσιμότητας στις τιμές και τον εντοπισμό περιπτώσεων συμφόρησης που οδηγούν σε διαφοροποίηση τιμών μεταξύ ζωνών. Τα αποτελέσματα περιλαμβάνουν όχι μόνο τις τιμές ανά ζώνη αλλά και τις ροές μεταφοράς μεταξύ ζωνών.

\subsection{Ανταγωνιστικό Αντιπαράδειγμα με Επικυρωμένη Ακρίβεια}
\label{subsec:counterfactual-contribution}

Σημαντική συνεισφορά αποτελεί η μεθοδολογία του βιβλίου εντολών οριακού κόστους (Κεφάλαιο 4), η οποία κατασκευάζει ένα \emph{ανταγωνιστικό αντιπαράδειγμα} της αγοράς: τις τιμές που θα προέκυπταν εάν κάθε συμμετέχων προσέφερε στο βραχυχρόνιο οριακό του κόστος συν την ορθολογική διαχείριση των τεχνικών του περιορισμών. Η προσέγγιση επιτυγχάνει συσχέτιση 0{,}84 σε ανάλυση δεκαπενταλέπτου χωρίς καμία στατιστική προσαρμογή εκ των υστέρων, διατηρώντας πλήρη ερμηνευσιμότητα: κάθε υπολογισμένη τιμή αντιστοιχεί σε συγκεκριμένη οριακή εντολή συγκεκριμένης τεχνολογίας. Πέραν της προσομοίωσης, η ιδιότητα αυτή καθιστά το σύστημα εργαλείο εποπτείας αγοράς: συστηματικές αποκλίσεις των πραγματικών τιμών \emph{πάνω} από το αντιπαράδειγμα, ιδίως όταν το ενδοημερήσιο σχήμα εξηγείται πλήρως, συνιστούν ποσοτικό μέτρο της απόκλισης της πραγματικής συμπεριφοράς προσφορών από την ανταγωνιστική (Ενότητα 5.7).

\subsection{Επίδειξη του Οικοσυστήματος Julia/JuMP}
\label{subsec:julia-contribution}

Η τέταρτη συνεισφορά είναι η επίδειξη της καταλληλότητας του οικοσυστήματος Julia/JuMP για προβλήματα βελτιστοποίησης στον τομέα της ενέργειας. Η Julia συνδυάζει την εκφραστικότητα γλωσσών υψηλού επιπέδου με την υπολογιστική απόδοση γλωσσών χαμηλού επιπέδου, ενώ το JuMP παρέχει μια διαισθητική διεπαφή για τη μοντελοποίηση προβλημάτων βελτιστοποίησης.

Η εργασία επιδεικνύει τη χρήση του JuMP για τη διατύπωση πολύπλοκων προβλημάτων βελτιστοποίησης, τη σύνδεση με διάφορους επιλυτές (HiGHS, Gurobi), και την αλληλεπίδραση με βάσεις δεδομένων. Ο κώδικας μπορεί να χρησιμεύσει ως παράδειγμα για παρόμοιες εφαρμογές.

\subsection{Μοντελοποίηση Ειδικών Παραμέτρων ανά Τύπο Καυσίμου}
\label{subsec:fuel-type-contribution}

Η πέμπτη συνεισφορά είναι η λεπτομερής μοντελοποίηση των τεχνικών χαρακτηριστικών των μονάδων παραγωγής ανάλογα με τον τύπο καυσίμου. Το σύστημα διαφοροποιεί τους χρόνους εκκίνησης, τους ρυθμούς μεταβολής, και τους ελάχιστους χρόνους λειτουργίας μεταξύ διαφορετικών τύπων καυσίμου όπως λιγνίτης, φυσικό αέριο, και υδροηλεκτρικά.

Η διαφοροποίηση αυτή οδηγεί σε ρεαλιστικότερα αποτελέσματα Unit Commitment και κατ' επέκταση σε ακριβέστερες προσφορές στην αγορά.


% =============================================================================
\section{Περιορισμοί}
\label{sec:limitations}

Παρά τις συνεισφορές που αναφέρθηκαν, η παρούσα υλοποίηση υπόκειται σε ορισμένους περιορισμούς που πρέπει να ληφθούν υπόψη κατά την ερμηνεία των αποτελεσμάτων.

\subsection{Μονοπερίοδη Βελτιστοποίηση}
\label{subsec:single-period-limitation}

Το σύστημα εκτελεί βελτιστοποίηση για μία ημέρα κάθε φορά, χωρίς να λαμβάνει υπόψη την αλληλεξάρτηση μεταξύ διαδοχικών ημερών. Στην πραγματικότητα, οι αποφάσεις δέσμευσης μονάδων επηρεάζονται από τις συνθήκες των προηγούμενων ημερών (αρχικές καταστάσεις) και τις προβλέψεις για τις επόμενες ημέρες (κυλιόμενος ορίζοντας).

Η απουσία κυλιόμενου ορίζοντα μπορεί να οδηγήσει σε μη βέλτιστες αποφάσεις, ιδιαίτερα για μονάδες με μεγάλους χρόνους εκκίνησης που πρέπει να ξεκινήσουν νωρίς για να είναι διαθέσιμες σε περιόδους υψηλής ζήτησης.

\subsection{Απλοποιημένο Μοντέλο Δικτύου}
\label{subsec:network-limitation}

Το μοντέλο δικτύου βασίζεται αποκλειστικά στη διαθέσιμη ικανότητα μεταφοράς (ATC) μεταξύ ζωνών προσφοράς, χωρίς να μοντελοποιεί τους νόμους της ροής ισχύος (Kirchhoff). Η προσέγγιση ATC είναι η τυπική για τις αγορές ηλεκτρικής ενέργειας, αλλά δεν αποτυπώνει με ακρίβεια τη φυσική συμπεριφορά του δικτύου.

Επιπλέον, δεν μοντελοποιούνται οι απώλειες μεταφοράς, οι οποίες μπορεί να είναι σημαντικές για μεγάλες αποστάσεις μεταφοράς.

\subsection{Περιορισμένοι Τύποι Εντολών}
\label{subsec:order-types-limitation}

Η υλοποίηση υποστηρίζει κυρίως απλές ωριαίες εντολές (Simple Orders), ενώ ο πραγματικός αλγόριθμος EUPHEMIA υποστηρίζει πολύ πιο σύνθετους τύπους εντολών. Οι εντολές μπλοκ (Block Orders) που επιτρέπουν την υποβολή προσφοράς για πολλαπλές συνεχόμενες ώρες υπό συνθήκη ελάχιστου εσόδου δεν υλοποιούνται πλήρως.

Επίσης, δεν υποστηρίζονται οι συνδεδεμένες εντολές μπλοκ (Linked Block Orders), οι αποκλειστικές εντολές μπλοκ (Exclusive Block Orders), οι εντολές ελάχιστου εισοδήματος (MIC Orders), και άλλοι προηγμένοι τύποι.

\subsection{Εκτίμηση Οικονομικών Παραμέτρων}
\label{subsec:economic-parameters-limitation}

Τα οικονομικά χαρακτηριστικά των μονάδων παραγωγής (οριακά κόστη, κόστη εκκίνησης, κόστη no-load) δεν δημοσιεύονται από το ENTSO-E και πρέπει να εκτιμηθούν με βάση τον τύπο καυσίμου και δημόσια διαθέσιμες πληροφορίες για τις τιμές καυσίμων και τους ρύπους CO$_2$.

Η αβεβαιότητα στην εκτίμηση αυτών των παραμέτρων αποτελεί σημαντική πηγή σφάλματος στα αποτελέσματα της προσομοίωσης.

\subsection{Μη Μοντελοποίηση Στρατηγικής Συμπεριφοράς}
\label{subsec:strategic-behavior-limitation}

Το σύστημα υποθέτει ότι οι συμμετέχοντες στην αγορά υποβάλλουν προσφορές με βάση το πραγματικό τους κόστος (ή με σταθερό συντελεστή markup). Στην πραγματικότητα, οι συμμετέχοντες μπορεί να ακολουθούν στρατηγικές υποβολής προσφορών που λαμβάνουν υπόψη τη θέση τους στην αγορά και τις αναμενόμενες ενέργειες των ανταγωνιστών.

Η επιλογή αυτή είναι εν μέρει σκόπιμη: καθιστά το σύστημα ανταγωνιστικό αντιπαράδειγμα, οπότε οι αποκλίσεις των πραγματικών τιμών από τις υπολογισμένες δεν αποτελούν αποκλειστικά σφάλμα μοντελοποίησης αλλά και μετρήσιμη ένδειξη μη ανταγωνιστικής συμπεριφοράς προσφορών. Περιορίζει ωστόσο την ικανότητα αναπαραγωγής των πραγματικών τιμών σε ημέρες στενότητας, όπου η στρατηγική συμπεριφορά είναι εντονότερη.


% =============================================================================
\section{Μελλοντική Εργασία}
\label{sec:future-work}

Με βάση την εμπειρία που αποκτήθηκε κατά την ανάπτυξη του συστήματος και τους περιορισμούς που εντοπίστηκαν, προτείνονται οι ακόλουθες κατευθύνσεις για μελλοντική έρευνα και ανάπτυξη.

\subsection{Συστηματοποίηση της Εποπτείας Αγοράς}
\label{subsec:market-monitoring-future}

Άμεση προέκταση της παρούσας εργασίας είναι η συστηματοποίηση της ανάλυσης καταλοίπων ως μεθοδολογίας εποπτείας αγοράς. Η ανάλυση σε επίπεδο μονάδας που παρουσιάστηκε ενδεικτικά στο Κεφάλαιο 5 (σύγκριση πραγματικής κατανομής με την αναμενόμενη από το ανταγωνιστικό αντιπαράδειγμα, εντοπισμός αδήλωτα μη διαθέσιμης ισχύος, δείκτης εναπομένοντος προμηθευτή) μπορεί να αυτοματοποιηθεί σε όλες τις ζώνες και ημέρες, με παλινδρόμηση των συναγόμενων περιθωρίων προσφορών σε εκ των προτέρων παρατηρήσιμα σήματα (περιθώριο επάρκειας, πρόβλεψη ΑΠΕ, πληρότητα ταμιευτήρων, ανακοινωθείσες μη διαθεσιμότητες) και έλεγχο συσχέτισης μεταξύ εταιρειών πέραν των κοινών σημάτων.

\subsection{Επέκταση στη Σύζευξη Βάσει Ροών}
\label{subsec:flow-based-future}

Η ανάλυση ανέδειξε ότι μετά τον Ιούνιο του 2022 τα σύνορα της περιοχής Core εκκαθαρίζονται με σύζευξη βάσει ροών (flow-based market coupling) και δεν αναπαρίστανται από δεδομένα ATC. Η υλοποίηση της flow-based προσέγγισης (πολύγωνα εφικτότητας από PTDF και περιθώρια κρίσιμων στοιχείων δικτύου) θα επέτρεπε την ενδογενή αναπαράσταση ζωνών όπως η Ουγγαρία, οι οποίες σήμερα αποδίδονται ατελώς ως δέκτες τιμών.

\subsection{Επέκταση σε Ενδοημερήσια Αγορά}
\label{subsec:intraday-extension}

Η πρώτη κατεύθυνση είναι η επέκταση του συστήματος για την προσομοίωση της Ενδοημερήσιας Αγοράς (Intraday Market). Η Ενδοημερήσια Αγορά λειτουργεί συνεχώς μεταξύ της λήξης της Προημερήσιας και της έναρξης της ημέρας παράδοσης, επιτρέποντας στους συμμετέχοντες να προσαρμόζουν τις θέσεις τους με βάση ενημερωμένες προβλέψεις.

Η μοντελοποίηση της Ενδοημερήσιας Αγοράς απαιτεί διαφορετικούς μηχανισμούς εκκαθάρισης (συνεχής διαπραγμάτευση vs δημοπρασία) και θα παρείχε πληρέστερη εικόνα της λειτουργίας της αγοράς.

\subsection{Προσομοίωση Αγοράς Εξισορρόπησης}
\label{subsec:balancing-simulation}

Η δεύτερη κατεύθυνση είναι η προσομοίωση της Αγοράς Εξισορρόπησης (Balancing Market). Η αγορά αυτή λειτουργεί σε πραγματικό χρόνο για την αντιμετώπιση αποκλίσεων μεταξύ παραγωγής και κατανάλωσης, και αποτελεί κρίσιμο συστατικό της λειτουργίας του συστήματος ηλεκτρικής ενέργειας.

Η μοντελοποίηση της Αγοράς Εξισορρόπησης θα απαιτήσει τη συμπερίληψη αβεβαιοτήτων στις προβλέψεις ζήτησης και παραγωγής ΑΠΕ, καθώς και τη μοντελοποίηση των αποθεμάτων ισχύος (reserves).

\subsection{Ενσωμάτωση Μηχανικής Μάθησης}
\label{subsec:ml-integration}

Η τρίτη κατεύθυνση είναι η ενσωμάτωση τεχνικών μηχανικής μάθησης για τη βελτίωση διαφόρων συνιστωσών του συστήματος. Συγκεκριμένα, η μηχανική μάθηση μπορεί να χρησιμοποιηθεί για την πρόβλεψη τιμών ηλεκτρικής ενέργειας, την εκτίμηση οικονομικών παραμέτρων μονάδων παραγωγής, και την πρόβλεψη διαθεσιμότητας μονάδων βάσει ιστορικών δεδομένων μη διαθεσιμότητας.

Η χρήση υβριδικών μοντέλων που συνδυάζουν τη φυσική γνώση (μοντέλα βελτιστοποίησης) με τη μηχανική μάθηση (data-driven approaches) αποτελεί πολλά υποσχόμενη κατεύθυνση.

\subsection{Υλοποίηση Σύνθετων Τύπων Εντολών}
\label{subsec:complex-orders}

Η τέταρτη κατεύθυνση είναι η πλήρης υλοποίηση των σύνθετων τύπων εντολών που υποστηρίζει ο αλγόριθμος EUPHEMIA. Οι εντολές μπλοκ με συνθήκη ελάχιστου εσόδου, οι συνδεδεμένες εντολές, και οι εντολές ελάχιστου εισοδήματος είναι απαραίτητες για τη ρεαλιστική αναπαράσταση της συμπεριφοράς θερμικών μονάδων στην αγορά.

Η υλοποίηση αυτών των τύπων εντολών θα απαιτήσει την επέκταση του μοντέλου MPCC με επιπλέον μεταβλητές και περιορισμούς.

\subsection{Βελτιστοποίηση Υπολογιστικής Απόδοσης}
\label{subsec:performance-optimization}

Η πέμπτη κατεύθυνση είναι η βελτιστοποίηση της υπολογιστικής απόδοσης του συστήματος για τη διαχείριση μεγαλύτερων προβλημάτων. Η παραλληλοποίηση της επίλυσης για πολλαπλές ημερομηνίες, η χρήση τεχνικών warm-start για διαδοχικές επιλύσεις, και η εξερεύνηση εναλλακτικών διατυπώσεων του προβλήματος αποτελούν πιθανές κατευθύνσεις.

Επιπλέον, η ενσωμάτωση του συστήματος σε περιβάλλον cloud θα επέτρεπε την κλιμάκωση σε μεγαλύτερα προβλήματα και την εκτέλεση εκτεταμένων αναλύσεων ευαισθησίας.

\subsection{Επέκταση σε Περιφερειακές Αγορές}
\label{subsec:regional-extension}

Η έκτη κατεύθυνση είναι η επέκταση της ανάλυσης σε περισσότερες ζώνες προσφοράς και η μελέτη περιφερειακών αγορών πέρα από τη Νοτιοανατολική Ευρώπη. Η σύγκριση των αποτελεσμάτων μεταξύ διαφορετικών περιφερειών θα παρείχε πληροφορίες για τις διαφορές στη δομή και λειτουργία των αγορών.

Η επέκταση αυτή θα απαιτήσει την αντιμετώπιση διαφορών στη διαθεσιμότητα και ποιότητα δεδομένων μεταξύ ζωνών.


% =============================================================================
\section{Επίλογος}
\label{sec:epilogue}

Η παρούσα διπλωματική εργασία επιχείρησε να συμβάλει στην κατανόηση και μοντελοποίηση των αγορών ηλεκτρικής ενέργειας μέσω της ανάπτυξης ενός ολοκληρωμένου συστήματος προσομοίωσης. Η εργασία κινείται στη διασταύρωση πολλών επιστημονικών πεδίων: της επιχειρησιακής έρευνας και βελτιστοποίησης, της οικονομικής θεωρίας των αγορών, της μηχανικής δεδομένων, και της ανάπτυξης λογισμικού.

Η σημασία της κατανόησης των μηχανισμών της αγοράς ηλεκτρικής ενέργειας αυξάνεται συνεχώς καθώς η ενεργειακή μετάβαση επιταχύνεται. Η διείσδυση των ανανεώσιμων πηγών ενέργειας, η ηλεκτροκίνηση, και η αποκέντρωση της παραγωγής δημιουργούν νέες προκλήσεις για τη λειτουργία των αγορών. Τα εργαλεία προσομοίωσης όπως αυτό που αναπτύχθηκε στην παρούσα εργασία μπορούν να συμβάλουν στην ανάλυση των επιπτώσεων αυτών των αλλαγών και στον σχεδιασμό πολιτικών.

Η επιλογή της Julia και του JuMP για την υλοποίηση αποδείχθηκε επιτυχής, επιτρέποντας τη διατύπωση των προβλημάτων βελτιστοποίησης με τρόπο που μοιάζει με τη μαθηματική διατύπωση, διατηρώντας παράλληλα υψηλή υπολογιστική απόδοση. Η εμπειρία αυτή επιβεβαιώνει ότι το οικοσύστημα Julia είναι ώριμο για χρήση σε εφαρμογές βελτιστοποίησης στον τομέα της ενέργειας.

Κλείνοντας, η εργασία αυτή αποτελεί μια βάση πάνω στην οποία μπορεί να χτιστεί περαιτέρω έρευνα και ανάπτυξη. Η διάθεση του κώδικα ως ανοιχτού λογισμικού έχει ως στόχο να διευκολύνει αυτή τη διαδικασία και να συμβάλει στην ευρύτερη κοινότητα που ασχολείται με τη μοντελοποίηση ενεργειακών συστημάτων.

% =============================================================================

\chapter{Συμπεράσματα}
\label{ch:conclusions}

Στο παρόν κεφάλαιο συνοψίζονται τα κύρια ευρήματα και οι συνεισφορές της διπλωματικής εργασίας, αναλύονται οι περιορισμοί της παρούσας υλοποίησης, και προτείνονται κατευθύνσεις για μελλοντική έρευνα και ανάπτυξη.


% =============================================================================
\section{Σύνοψη Εργασίας}
\label{sec:thesis-summary}

Η παρούσα διπλωματική εργασία ασχολήθηκε με τη μοντελοποίηση και προσομοίωση της Προημερήσιας Αγοράς ηλεκτρικής ενέργειας, με έμφαση στον αλγόριθμο EUPHEMIA που χρησιμοποιείται για την εκκαθάριση των ευρωπαϊκών αγορών. Αναπτύχθηκε ένα ολοκληρωμένο σύστημα λογισμικού που καλύπτει όλα τα στάδια της διαδικασίας, από τη συλλογή δεδομένων έως τον υπολογισμό τιμών εκκαθάρισης.

Η εργασία ξεκίνησε με την ανάλυση του πλαισίου λειτουργίας των αγορών ηλεκτρικής ενέργειας στην Ευρώπη, εξετάζοντας την ιστορική εξέλιξη από τα κρατικά μονοπώλια στις απελευθερωμένες αγορές. Παρουσιάστηκε η δομή της Προημερήσιας Αγοράς και ο ρόλος του αλγορίθμου EUPHEMIA στην ενοποίηση των ευρωπαϊκών αγορών μέσω του μηχανισμού Single Day-Ahead Coupling (SDAC).

Στη συνέχεια, αναπτύχθηκε το θεωρητικό υπόβαθρο που απαιτείται για την κατανόηση του προβλήματος. Παρουσιάστηκαν οι βασικές έννοιες του μαθηματικού προγραμματισμού, με έμφαση στα προβλήματα μεικτού ακέραιου προγραμματισμού (MIP) και στα προβλήματα συμπληρωματικότητας (MPCC). Διατυπώθηκε το πρόβλημα Unit Commitment και αναλύθηκε ο αλγόριθμος EUPHEMIA με βάση τη δημόσια τεκμηρίωση.

Η αρχιτεκτονική του συστήματος σχεδιάστηκε με γνώμονα την επεκτασιμότητα και τη συντηρησιμότητα. Αναπτύχθηκε αγωγός δεδομένων (data pipeline) σε Python για τη συλλογή δεδομένων από την πλατφόρμα διαφάνειας του ENTSO-E, με αποθήκευση σε βάση δεδομένων PostgreSQL. Η κύρια εφαρμογή υλοποιήθηκε στη γλώσσα Julia με χρήση του πλαισίου JuMP για τη μοντελοποίηση των προβλημάτων βελτιστοποίησης.

Η υλοποίηση περιλαμβάνει τρία κύρια συστατικά. Πρώτον, τον επιλυτή Unit Commitment που βελτιστοποιεί τη δέσμευση και την κατανομή των μονάδων παραγωγής λαμβάνοντας υπόψη τους τεχνικούς περιορισμούς κάθε τύπου καυσίμου. Δεύτερον, τη μηχανή στρατηγικής προσφορών που μετατρέπει τα αποτελέσματα του Unit Commitment σε εντολές αγοράς με διάφορες στρατηγικές. Τρίτον, τον αλγόριθμο εκκαθάρισης αγοράς βασισμένο στη μέθοδο MPCC, που υπολογίζει τις τιμές εκκαθάρισης και τους βαθμούς αποδοχής των εντολών.

Τέλος, το σύστημα αξιολογήθηκε μέσω σύγκρισης των υπολογισμένων τιμών με τις πραγματικές τιμές της αγοράς. Το βιβλίο εντολών οριακού κόστους, τροφοδοτούμενο με τους κύριους θεμελιώδεις παράγοντες της αγοράς (ημερήσιες τιμές φυσικού αερίου TTF και δικαιωμάτων CO$_2$, αξία νερού συναρτήσει της πληρότητας των ταμιευτήρων, αυτοπρογραμματισμός μονάδων βάσης, παρατηρούμενες καθαρές εισαγωγές), αναπαράγει τις τιμές της ελληνικής ζώνης με συντελεστή συσχέτισης $r = 0{,}84$ σε ανάλυση δεκαπενταλέπτου επί συνεχούς τετραμήνου, με πρακτικά μηδενική συστηματική απόκλιση και διάμεσο ενδοημερήσιο συντελεστή συσχέτισης 0{,}90 (Κεφάλαιο 5).


% =============================================================================
\section{Κύριες Συνεισφορές}
\label{sec:contributions}

Η παρούσα εργασία παρουσιάζει τις ακόλουθες συνεισφορές στους τομείς της μοντελοποίησης αγορών ενέργειας και της εφαρμοσμένης βελτιστοποίησης.

\subsection{Υλοποίηση Ανοιχτού Κώδικα}
\label{subsec:open-source-contribution}

Η πρώτη και κυριότερη συνεισφορά είναι η ανάπτυξη μιας πλήρως λειτουργικής υλοποίησης ανοιχτού κώδικα για την προσομοίωση της Προημερήσιας Αγοράς. Ενώ ο αλγόριθμος EUPHEMIA είναι δημόσια τεκμηριωμένος, οι υπάρχουσες υλοποιήσεις είναι είτε ιδιόκτητες είτε περιορισμένης λειτουργικότητας. Εξ όσων είμαστε σε θέση να γνωρίζουμε, πρόκειται για το πρώτο σύστημα ανοιχτού κώδικα που μοντελοποιεί την αγορά αυτή με τεκμηριωμένη, ποσοτικά επικυρωμένη ακρίβεια σε πραγματικά δεδομένα και σε συνεχή περίοδο λειτουργίας, με πλήρως αναπαραγώγιμη μεθοδολογία --- από τα ακατέργαστα δημόσια δεδομένα έως τις τιμές εκκαθάρισης. Το σύστημα που αναπτύχθηκε στην παρούσα εργασία διατίθεται ελεύθερα στη διεύθυνση \url{https://github.com/philokalia-ai/energy-markets} και μπορεί να χρησιμοποιηθεί για εκπαιδευτικούς και ερευνητικούς σκοπούς.

Η υλοποίηση περιλαμβάνει όλα τα απαραίτητα συστατικά για την εκτέλεση προσομοιώσεων, από τη συλλογή δεδομένων έως την αποθήκευση αποτελεσμάτων. Ο κώδικας είναι καλά τεκμηριωμένος και δομημένος με τρόπο που διευκολύνει την επέκταση και την προσαρμογή σε διαφορετικές ανάγκες.

\subsection{Ολοκληρωμένος Αγωγός Δεδομένων}
\label{subsec:data-pipeline-contribution}

Η δεύτερη συνεισφορά είναι η ανάπτυξη ενός πλήρους αγωγού δεδομένων για τη συλλογή και επεξεργασία δεδομένων από την πλατφόρμα διαφάνειας του ENTSO-E. Ο αγωγός υποστηρίζει σταδιακές ενημερώσεις (incremental updates), αποτελεσματική μαζική φόρτωση δεδομένων, και αυτοματοποιημένη εκτέλεση μέσω GitHub Actions.

Τα δεδομένα που συλλέγονται περιλαμβάνουν πληροφορίες για τις μονάδες παραγωγής, το ιστορικό φορτίο, τις προβλέψεις παραγωγής από ΑΠΕ, τις διαθέσιμες ικανότητες μεταφοράς, και τις μη διαθεσιμότητες μονάδων. Ο σχεδιασμός του σχήματος της βάσης δεδομένων επιτρέπει την αποτελεσματική εκτέλεση ερωτημάτων και τη διατήρηση της ιστορικότητας των δεδομένων.

\subsection{Πολυζωνική Εκκαθάριση Αγοράς}
\label{subsec:multizone-contribution}

Η τρίτη συνεισφορά είναι η υλοποίηση πολυζωνικής εκκαθάρισης αγοράς με περιορισμούς διασυνοριακής μεταφοράς. Το σύστημα μοντελοποιεί τη διαθέσιμη ικανότητα μεταφοράς (Available Transfer Capacity - ATC) μεταξύ ζωνών προσφοράς και υπολογίζει ταυτόχρονα τις τιμές σε όλες τις συνδεδεμένες ζώνες.

Η προσέγγιση αυτή επιτρέπει την ανάλυση της επίδρασης της διασυνδεσιμότητας στις τιμές και τον εντοπισμό περιπτώσεων συμφόρησης που οδηγούν σε διαφοροποίηση τιμών μεταξύ ζωνών. Τα αποτελέσματα περιλαμβάνουν όχι μόνο τις τιμές ανά ζώνη αλλά και τις ροές μεταφοράς μεταξύ ζωνών.

\subsection{Ανταγωνιστικό Αντιπαράδειγμα με Επικυρωμένη Ακρίβεια}
\label{subsec:counterfactual-contribution}

Σημαντική συνεισφορά αποτελεί η μεθοδολογία του βιβλίου εντολών οριακού κόστους (Κεφάλαιο 4), η οποία κατασκευάζει ένα \emph{ανταγωνιστικό αντιπαράδειγμα} της αγοράς: τις τιμές που θα προέκυπταν εάν κάθε συμμετέχων προσέφερε στο βραχυχρόνιο οριακό του κόστος συν την ορθολογική διαχείριση των τεχνικών του περιορισμών. Η προσέγγιση επιτυγχάνει συσχέτιση 0{,}84 σε ανάλυση δεκαπενταλέπτου χωρίς καμία στατιστική προσαρμογή εκ των υστέρων, διατηρώντας πλήρη ερμηνευσιμότητα: κάθε υπολογισμένη τιμή αντιστοιχεί σε συγκεκριμένη οριακή εντολή συγκεκριμένης τεχνολογίας. Πέραν της προσομοίωσης, η ιδιότητα αυτή καθιστά το σύστημα εργαλείο εποπτείας αγοράς: συστηματικές αποκλίσεις των πραγματικών τιμών \emph{πάνω} από το αντιπαράδειγμα, ιδίως όταν το ενδοημερήσιο σχήμα εξηγείται πλήρως, συνιστούν ποσοτικό μέτρο της απόκλισης της πραγματικής συμπεριφοράς προσφορών από την ανταγωνιστική (Ενότητα 5.7).

\subsection{Επίδειξη του Οικοσυστήματος Julia/JuMP}
\label{subsec:julia-contribution}

Η τέταρτη συνεισφορά είναι η επίδειξη της καταλληλότητας του οικοσυστήματος Julia/JuMP για προβλήματα βελτιστοποίησης στον τομέα της ενέργειας. Η Julia συνδυάζει την εκφραστικότητα γλωσσών υψηλού επιπέδου με την υπολογιστική απόδοση γλωσσών χαμηλού επιπέδου, ενώ το JuMP παρέχει μια διαισθητική διεπαφή για τη μοντελοποίηση προβλημάτων βελτιστοποίησης.

Η εργασία επιδεικνύει τη χρήση του JuMP για τη διατύπωση πολύπλοκων προβλημάτων βελτιστοποίησης, τη σύνδεση με διάφορους επιλυτές (HiGHS, Gurobi), και την αλληλεπίδραση με βάσεις δεδομένων. Ο κώδικας μπορεί να χρησιμεύσει ως παράδειγμα για παρόμοιες εφαρμογές.

\subsection{Μοντελοποίηση Ειδικών Παραμέτρων ανά Τύπο Καυσίμου}
\label{subsec:fuel-type-contribution}

Η πέμπτη συνεισφορά είναι η λεπτομερής μοντελοποίηση των τεχνικών χαρακτηριστικών των μονάδων παραγωγής ανάλογα με τον τύπο καυσίμου. Το σύστημα διαφοροποιεί τους χρόνους εκκίνησης, τους ρυθμούς μεταβολής, και τους ελάχιστους χρόνους λειτουργίας μεταξύ διαφορετικών τύπων καυσίμου όπως λιγνίτης, φυσικό αέριο, και υδροηλεκτρικά.

Η διαφοροποίηση αυτή οδηγεί σε ρεαλιστικότερα αποτελέσματα Unit Commitment και κατ' επέκταση σε ακριβέστερες προσφορές στην αγορά.


% =============================================================================
\section{Περιορισμοί}
\label{sec:limitations}

Παρά τις συνεισφορές που αναφέρθηκαν, η παρούσα υλοποίηση υπόκειται σε ορισμένους περιορισμούς που πρέπει να ληφθούν υπόψη κατά την ερμηνεία των αποτελεσμάτων.

\subsection{Μονοπερίοδη Βελτιστοποίηση}
\label{subsec:single-period-limitation}

Το σύστημα εκτελεί βελτιστοποίηση για μία ημέρα κάθε φορά, χωρίς να λαμβάνει υπόψη την αλληλεξάρτηση μεταξύ διαδοχικών ημερών. Στην πραγματικότητα, οι αποφάσεις δέσμευσης μονάδων επηρεάζονται από τις συνθήκες των προηγούμενων ημερών (αρχικές καταστάσεις) και τις προβλέψεις για τις επόμενες ημέρες (κυλιόμενος ορίζοντας).

Η απουσία κυλιόμενου ορίζοντα μπορεί να οδηγήσει σε μη βέλτιστες αποφάσεις, ιδιαίτερα για μονάδες με μεγάλους χρόνους εκκίνησης που πρέπει να ξεκινήσουν νωρίς για να είναι διαθέσιμες σε περιόδους υψηλής ζήτησης.

\subsection{Απλοποιημένο Μοντέλο Δικτύου}
\label{subsec:network-limitation}

Το μοντέλο δικτύου βασίζεται αποκλειστικά στη διαθέσιμη ικανότητα μεταφοράς (ATC) μεταξύ ζωνών προσφοράς, χωρίς να μοντελοποιεί τους νόμους της ροής ισχύος (Kirchhoff). Η προσέγγιση ATC είναι η τυπική για τις αγορές ηλεκτρικής ενέργειας, αλλά δεν αποτυπώνει με ακρίβεια τη φυσική συμπεριφορά του δικτύου.

Επιπλέον, δεν μοντελοποιούνται οι απώλειες μεταφοράς, οι οποίες μπορεί να είναι σημαντικές για μεγάλες αποστάσεις μεταφοράς.

\subsection{Περιορισμένοι Τύποι Εντολών}
\label{subsec:order-types-limitation}

Η υλοποίηση υποστηρίζει κυρίως απλές ωριαίες εντολές (Simple Orders), ενώ ο πραγματικός αλγόριθμος EUPHEMIA υποστηρίζει πολύ πιο σύνθετους τύπους εντολών. Οι εντολές μπλοκ (Block Orders) που επιτρέπουν την υποβολή προσφοράς για πολλαπλές συνεχόμενες ώρες υπό συνθήκη ελάχιστου εσόδου δεν υλοποιούνται πλήρως.

Επίσης, δεν υποστηρίζονται οι συνδεδεμένες εντολές μπλοκ (Linked Block Orders), οι αποκλειστικές εντολές μπλοκ (Exclusive Block Orders), οι εντολές ελάχιστου εισοδήματος (MIC Orders), και άλλοι προηγμένοι τύποι.

\subsection{Εκτίμηση Οικονομικών Παραμέτρων}
\label{subsec:economic-parameters-limitation}

Τα οικονομικά χαρακτηριστικά των μονάδων παραγωγής (οριακά κόστη, κόστη εκκίνησης, κόστη no-load) δεν δημοσιεύονται από το ENTSO-E και πρέπει να εκτιμηθούν με βάση τον τύπο καυσίμου και δημόσια διαθέσιμες πληροφορίες για τις τιμές καυσίμων και τους ρύπους CO$_2$.

Η αβεβαιότητα στην εκτίμηση αυτών των παραμέτρων αποτελεί σημαντική πηγή σφάλματος στα αποτελέσματα της προσομοίωσης.

\subsection{Μη Μοντελοποίηση Στρατηγικής Συμπεριφοράς}
\label{subsec:strategic-behavior-limitation}

Το σύστημα υποθέτει ότι οι συμμετέχοντες στην αγορά υποβάλλουν προσφορές με βάση το πραγματικό τους κόστος (ή με σταθερό συντελεστή markup). Στην πραγματικότητα, οι συμμετέχοντες μπορεί να ακολουθούν στρατηγικές υποβολής προσφορών που λαμβάνουν υπόψη τη θέση τους στην αγορά και τις αναμενόμενες ενέργειες των ανταγωνιστών.

Η επιλογή αυτή είναι εν μέρει σκόπιμη: καθιστά το σύστημα ανταγωνιστικό αντιπαράδειγμα, οπότε οι αποκλίσεις των πραγματικών τιμών από τις υπολογισμένες δεν αποτελούν αποκλειστικά σφάλμα μοντελοποίησης αλλά και μετρήσιμη ένδειξη μη ανταγωνιστικής συμπεριφοράς προσφορών. Περιορίζει ωστόσο την ικανότητα αναπαραγωγής των πραγματικών τιμών σε ημέρες στενότητας, όπου η στρατηγική συμπεριφορά είναι εντονότερη.


% =============================================================================
\section{Μελλοντική Εργασία}
\label{sec:future-work}

Με βάση την εμπειρία που αποκτήθηκε κατά την ανάπτυξη του συστήματος και τους περιορισμούς που εντοπίστηκαν, προτείνονται οι ακόλουθες κατευθύνσεις για μελλοντική έρευνα και ανάπτυξη.

\subsection{Συστηματοποίηση της Εποπτείας Αγοράς}
\label{subsec:market-monitoring-future}

Άμεση προέκταση της παρούσας εργασίας είναι η συστηματοποίηση της ανάλυσης καταλοίπων ως μεθοδολογίας εποπτείας αγοράς. Η ανάλυση σε επίπεδο μονάδας που παρουσιάστηκε ενδεικτικά στο Κεφάλαιο 5 (σύγκριση πραγματικής κατανομής με την αναμενόμενη από το ανταγωνιστικό αντιπαράδειγμα, εντοπισμός αδήλωτα μη διαθέσιμης ισχύος, δείκτης εναπομένοντος προμηθευτή) μπορεί να αυτοματοποιηθεί σε όλες τις ζώνες και ημέρες, με παλινδρόμηση των συναγόμενων περιθωρίων προσφορών σε εκ των προτέρων παρατηρήσιμα σήματα (περιθώριο επάρκειας, πρόβλεψη ΑΠΕ, πληρότητα ταμιευτήρων, ανακοινωθείσες μη διαθεσιμότητες) και έλεγχο συσχέτισης μεταξύ εταιρειών πέραν των κοινών σημάτων.

\subsection{Επέκταση στη Σύζευξη Βάσει Ροών}
\label{subsec:flow-based-future}

Η ανάλυση ανέδειξε ότι μετά τον Ιούνιο του 2022 τα σύνορα της περιοχής Core εκκαθαρίζονται με σύζευξη βάσει ροών (flow-based market coupling) και δεν αναπαρίστανται από δεδομένα ATC. Η υλοποίηση της flow-based προσέγγισης (πολύγωνα εφικτότητας από PTDF και περιθώρια κρίσιμων στοιχείων δικτύου) θα επέτρεπε την ενδογενή αναπαράσταση ζωνών όπως η Ουγγαρία, οι οποίες σήμερα αποδίδονται ατελώς ως δέκτες τιμών.

\subsection{Επέκταση σε Ενδοημερήσια Αγορά}
\label{subsec:intraday-extension}

Η πρώτη κατεύθυνση είναι η επέκταση του συστήματος για την προσομοίωση της Ενδοημερήσιας Αγοράς (Intraday Market). Η Ενδοημερήσια Αγορά λειτουργεί συνεχώς μεταξύ της λήξης της Προημερήσιας και της έναρξης της ημέρας παράδοσης, επιτρέποντας στους συμμετέχοντες να προσαρμόζουν τις θέσεις τους με βάση ενημερωμένες προβλέψεις.

Η μοντελοποίηση της Ενδοημερήσιας Αγοράς απαιτεί διαφορετικούς μηχανισμούς εκκαθάρισης (συνεχής διαπραγμάτευση vs δημοπρασία) και θα παρείχε πληρέστερη εικόνα της λειτουργίας της αγοράς.

\subsection{Προσομοίωση Αγοράς Εξισορρόπησης}
\label{subsec:balancing-simulation}

Η δεύτερη κατεύθυνση είναι η προσομοίωση της Αγοράς Εξισορρόπησης (Balancing Market). Η αγορά αυτή λειτουργεί σε πραγματικό χρόνο για την αντιμετώπιση αποκλίσεων μεταξύ παραγωγής και κατανάλωσης, και αποτελεί κρίσιμο συστατικό της λειτουργίας του συστήματος ηλεκτρικής ενέργειας.

Η μοντελοποίηση της Αγοράς Εξισορρόπησης θα απαιτήσει τη συμπερίληψη αβεβαιοτήτων στις προβλέψεις ζήτησης και παραγωγής ΑΠΕ, καθώς και τη μοντελοποίηση των αποθεμάτων ισχύος (reserves).

\subsection{Ενσωμάτωση Μηχανικής Μάθησης}
\label{subsec:ml-integration}

Η τρίτη κατεύθυνση είναι η ενσωμάτωση τεχνικών μηχανικής μάθησης για τη βελτίωση διαφόρων συνιστωσών του συστήματος. Συγκεκριμένα, η μηχανική μάθηση μπορεί να χρησιμοποιηθεί για την πρόβλεψη τιμών ηλεκτρικής ενέργειας, την εκτίμηση οικονομικών παραμέτρων μονάδων παραγωγής, και την πρόβλεψη διαθεσιμότητας μονάδων βάσει ιστορικών δεδομένων μη διαθεσιμότητας.

Η χρήση υβριδικών μοντέλων που συνδυάζουν τη φυσική γνώση (μοντέλα βελτιστοποίησης) με τη μηχανική μάθηση (data-driven approaches) αποτελεί πολλά υποσχόμενη κατεύθυνση.

\subsection{Υλοποίηση Σύνθετων Τύπων Εντολών}
\label{subsec:complex-orders}

Η τέταρτη κατεύθυνση είναι η πλήρης υλοποίηση των σύνθετων τύπων εντολών που υποστηρίζει ο αλγόριθμος EUPHEMIA. Οι εντολές μπλοκ με συνθήκη ελάχιστου εσόδου, οι συνδεδεμένες εντολές, και οι εντολές ελάχιστου εισοδήματος είναι απαραίτητες για τη ρεαλιστική αναπαράσταση της συμπεριφοράς θερμικών μονάδων στην αγορά.

Η υλοποίηση αυτών των τύπων εντολών θα απαιτήσει την επέκταση του μοντέλου MPCC με επιπλέον μεταβλητές και περιορισμούς.

\subsection{Βελτιστοποίηση Υπολογιστικής Απόδοσης}
\label{subsec:performance-optimization}

Η πέμπτη κατεύθυνση είναι η βελτιστοποίηση της υπολογιστικής απόδοσης του συστήματος για τη διαχείριση μεγαλύτερων προβλημάτων. Η παραλληλοποίηση της επίλυσης για πολλαπλές ημερομηνίες, η χρήση τεχνικών warm-start για διαδοχικές επιλύσεις, και η εξερεύνηση εναλλακτικών διατυπώσεων του προβλήματος αποτελούν πιθανές κατευθύνσεις.

Επιπλέον, η ενσωμάτωση του συστήματος σε περιβάλλον cloud θα επέτρεπε την κλιμάκωση σε μεγαλύτερα προβλήματα και την εκτέλεση εκτεταμένων αναλύσεων ευαισθησίας.

\subsection{Επέκταση σε Περιφερειακές Αγορές}
\label{subsec:regional-extension}

Η έκτη κατεύθυνση είναι η επέκταση της ανάλυσης σε περισσότερες ζώνες προσφοράς και η μελέτη περιφερειακών αγορών πέρα από τη Νοτιοανατολική Ευρώπη. Η σύγκριση των αποτελεσμάτων μεταξύ διαφορετικών περιφερειών θα παρείχε πληροφορίες για τις διαφορές στη δομή και λειτουργία των αγορών.

Η επέκταση αυτή θα απαιτήσει την αντιμετώπιση διαφορών στη διαθεσιμότητα και ποιότητα δεδομένων μεταξύ ζωνών.


% =============================================================================
\section{Επίλογος}
\label{sec:epilogue}

Η παρούσα διπλωματική εργασία επιχείρησε να συμβάλει στην κατανόηση και μοντελοποίηση των αγορών ηλεκτρικής ενέργειας μέσω της ανάπτυξης ενός ολοκληρωμένου συστήματος προσομοίωσης. Η εργασία κινείται στη διασταύρωση πολλών επιστημονικών πεδίων: της επιχειρησιακής έρευνας και βελτιστοποίησης, της οικονομικής θεωρίας των αγορών, της μηχανικής δεδομένων, και της ανάπτυξης λογισμικού.

Η σημασία της κατανόησης των μηχανισμών της αγοράς ηλεκτρικής ενέργειας αυξάνεται συνεχώς καθώς η ενεργειακή μετάβαση επιταχύνεται. Η διείσδυση των ανανεώσιμων πηγών ενέργειας, η ηλεκτροκίνηση, και η αποκέντρωση της παραγωγής δημιουργούν νέες προκλήσεις για τη λειτουργία των αγορών. Τα εργαλεία προσομοίωσης όπως αυτό που αναπτύχθηκε στην παρούσα εργασία μπορούν να συμβάλουν στην ανάλυση των επιπτώσεων αυτών των αλλαγών και στον σχεδιασμό πολιτικών.

Η επιλογή της Julia και του JuMP για την υλοποίηση αποδείχθηκε επιτυχής, επιτρέποντας τη διατύπωση των προβλημάτων βελτιστοποίησης με τρόπο που μοιάζει με τη μαθηματική διατύπωση, διατηρώντας παράλληλα υψηλή υπολογιστική απόδοση. Η εμπειρία αυτή επιβεβαιώνει ότι το οικοσύστημα Julia είναι ώριμο για χρήση σε εφαρμογές βελτιστοποίησης στον τομέα της ενέργειας.

Κλείνοντας, η εργασία αυτή αποτελεί μια βάση πάνω στην οποία μπορεί να χτιστεί περαιτέρω έρευνα και ανάπτυξη. Η διάθεση του κώδικα ως ανοιχτού λογισμικού έχει ως στόχο να διευκολύνει αυτή τη διαδικασία και να συμβάλει στην ευρύτερη κοινότητα που ασχολείται με τη μοντελοποίηση ενεργειακών συστημάτων.
